\documentclass[11pt,a4paper]{article}

% ---------------------------------------------------------------------------
% CJKQC user manual
% The update date is deliberately fixed. Change \ManualDate whenever the
% manual's content changes.
% ---------------------------------------------------------------------------
\newcommand{\ManualDate}{11 November 2026}
\newcommand{\ManualVersion}{1.0}

\usepackage[T1]{fontenc}
\usepackage[utf8]{inputenc}
\usepackage{lmodern}
\usepackage[a4paper,margin=2.5cm,headheight=14pt]{geometry}
\usepackage{microtype}
\usepackage{xcolor}
\usepackage{graphicx}
\usepackage{xstring}
\usepackage{amsmath}
\usepackage{amssymb}
\usepackage{array}
\usepackage{booktabs}
\usepackage{longtable}
\usepackage{tabularx}
\usepackage{enumitem}
\usepackage{fancyhdr}
\usepackage{titlesec}
\usepackage[breakable]{tcolorbox}
\usepackage{hyperref}

\definecolor{deepblue}{HTML}{12355B}
\definecolor{noteblue}{HTML}{E8F0F8}
\definecolor{warnorange}{HTML}{FFF1E0}

\hypersetup{
  colorlinks=true,
  linkcolor=deepblue,
  urlcolor=deepblue,
  citecolor=deepblue,
  pdftitle={CJKQC: A Vibey User Manual},
  pdfauthor={The CJKQC project},
  pdfsubject={Collecting, filtering, and thermochemical post-processing of quantum-chemistry output},
  pdfkeywords={CJKQC, quantum chemistry, molecular clusters, thermochemistry, JSON Lines}
}

\titleformat{\section}
  {\Large\bfseries\color{deepblue}}{\thesection}{0.6em}{}
\titleformat{\subsection}
  {\large\bfseries\color{deepblue}}{\thesubsection}{0.5em}{}
\titleformat{\subsubsection}
  {\normalsize\bfseries\color{deepblue}}{\thesubsubsection}{0.5em}{}

\setlist[itemize]{topsep=2pt,itemsep=1pt,parsep=1pt}
\setlist[enumerate]{topsep=2pt,itemsep=1pt,parsep=1pt}
\setlength{\emergencystretch}{2em}

% Ragged-right paragraph columns: long flag and field names cannot be
% hyphenated, so justified narrow columns stretch badly.
\newcolumntype{L}[1]{>{\raggedright\arraybackslash}p{#1}}

\pagestyle{fancy}
\fancyhf{}
\fancyhead[L]{\footnotesize CJKQC: A Vibey User Manual}
\fancyhead[R]{\footnotesize Updated \ManualDate}
\fancyfoot[C]{\thepage}
\renewcommand{\headrulewidth}{0.3pt}

% Convenience macros. \detokenize keeps underscores and other special
% characters in flags and JSON field names readable in the PDF.
\newcommand{\CJKQC}{\textsf{CJKQC}}

% A colourful, slightly shaken wordmark for the cover.
\newcommand{\vibey}{%
  \raisebox{-2pt}{\color{red!85!black}\rotatebox{-9}{V}}%
  \kern-1pt\raisebox{1pt}{\color{orange!90!black}\rotatebox{7}{i}}%
  \kern-1pt\raisebox{-1pt}{\color{magenta!80!black}\rotatebox{-6}{b}}%
  \kern-1pt\raisebox{2pt}{\color{green!60!black}\rotatebox{8}{e}}%
  \kern-1pt\raisebox{-2pt}{\color{blue!85!black}\rotatebox{-8}{y}}%
}
% \detokenize makes underscores and other characters printable. The
% substitution keeps two hyphens as two hyphens: T1 typewriter fonts
% otherwise turn "--" into an en dash.
\newcommand{\codefmt}[1]{%
  \texttt{\StrSubstitute{\detokenize{#1}}{--}{-\mbox{}-}}}
\newcommand{\flg}[1]{\codefmt{#1}}
\newcommand{\fld}[1]{\codefmt{#1}}
\newcommand{\val}[1]{\codefmt{#1}}
\newcommand{\unit}[1]{\,\mathrm{#1}}

\newtcolorbox{note}[1][]{
  breakable,
  colback=noteblue,
  colframe=deepblue,
  boxrule=0.4pt,
  left=6pt,right=6pt,top=4pt,bottom=4pt,
  fonttitle=\bfseries,
  title={Note},
  #1
}
\newtcolorbox{warning}[1][]{
  breakable,
  colback=warnorange,
  colframe=orange!70!black,
  boxrule=0.4pt,
  left=6pt,right=6pt,top=4pt,bottom=4pt,
  fonttitle=\bfseries,
  title={Warning},
  #1
}
\newtcolorbox{example}[1][]{
  breakable,
  colback=black!3,
  colframe=black!35,
  boxrule=0.4pt,
  left=6pt,right=6pt,top=4pt,bottom=4pt,
  fonttitle=\bfseries,
  title={Example},
  #1
}

\begin{document}

% Keep the hidden title page from reusing hyperref's page.1 anchor.
\pagenumbering{Alph}

% ---------------------------------------------------------------------------
% Title page
% ---------------------------------------------------------------------------
\begin{titlepage}
\thispagestyle{empty}
\centering
\vspace*{3.2cm}

{\color{deepblue}\rule{\textwidth}{1.2pt}}\\[0.8cm]
{\Huge\bfseries\color{deepblue} CJKQC\\[0.25cm] A \vibey{} User Manual\par}
\vspace{0.7cm}
{\Large Collecting, filtering, and thermochemical\\[0.15cm]
post-processing of quantum-chemistry output\par}
\vspace{0.6cm}
{\color{deepblue}\rule{\textwidth}{1.2pt}}\\[1.5cm]

{\large A practical guide to databases, printing, filtering,\\
thermochemistry, corrections, formation energies, and solvation\par}

\vfill

\begin{tabular}{rl}
\textbf{Covers:} & CJKQC 0.1.0\\[2pt]
\textbf{Manual version:} & \ManualVersion\\[2pt]
\textbf{Last updated:} & \ManualDate\\
\end{tabular}

\vspace{1.2cm}
{\footnotesize This manual describes the command-line tool \CJKQC{} and the data it reads and writes.\\
The requirements documents under \texttt{rqm/} remain the authoritative description of behavior.}
\end{titlepage}

% ---------------------------------------------------------------------------
% Front matter
% ---------------------------------------------------------------------------
\pagenumbering{roman}
\tableofcontents
\clearpage

\pagenumbering{arabic}

\section{Introduction and quick start}\label{sec:introduction}

\CJKQC{} is a command-line program for turning folders of
quantum-chemistry output into a compact database, and then for inspecting,
filtering, and post-processing that database. It is a spiritual successor to JKQC from
the JKCS2.1 framework. The practical workflow is always the same:

\begin{enumerate}
  \item \textbf{Collect} program outputs into a database file, with one
        record per line.
  \item \textbf{Inspect and filter} the records: print columns, select
        compositions, apply energy or frequency cutoffs, remove duplicates,
        and sort.
  \item \textbf{Post-process} the survivors: derive thermochemistry,
        attach single-point corrections, compute formation energies, and
        print solvation distributions.
\end{enumerate}

Nothing derived is stored. A database holds only what the quantum-chemistry
program actually printed (energies, geometries, frequencies, and so on).
Zero-point energies, Gibbs free energies, corrections, formation energies,
populations, and distributions are recomputed on demand from those
stored values under the settings of the current run. This is why changing the
temperature, a frequency scale factor, or an energy unit never invalidates a
database.

\subsection{What can be read}

\CJKQC{} chooses how to read each file by looking at its \emph{content}, never
at its name or extension. The supported formats are:

\begin{center}
\begin{tabular}{@{}ll@{}}
\toprule
Program & Typical input \\
\midrule
Gaussian & \texttt{.log}, \texttt{.out} files with a Gaussian banner \\
ORCA     & \texttt{.out} files with ORCA's text banner \\
xtb      & log files with an \texttt{xtb version} banner \\
MRCC     & output files containing \texttt{MRCC program system} \\
plain xyz & an atom count, a comment line, and atom rows \\
\bottomrule
\end{tabular}
\end{center}

One file becomes exactly one record. For a file that chains several jobs,
the record describes the \emph{final} job (the last SCF energy, the final
geometry and frequencies, the final property block). For a multi-frame xyz
file, it describes the last complete frame. Failed and truncated jobs are
stored as records with status \fld{failed} or \fld{truncated}; they do not
abort a collection run. A file that matches no supported format, or an empty file,
\emph{does} abort a build run, because it is not a quantum-chemistry job at
all.

\subsection{Command anatomy}

Every invocation has the same shape:

\begin{verbatim}
CJKQC <inputs>... [filtering flags] [print flags]
      [settings flags] [-o out.jsonl]
\end{verbatim}

The positional \texttt{<inputs>} are either quantum-chemistry output files or
existing databases. \CJKQC{} classifies them by content:

\begin{itemize}
  \item \textbf{all inputs are databases} \(\rightarrow\) \textbf{merge mode}:
        the databases are combined, and records are copied exactly as stored;
  \item \textbf{all inputs are program outputs} \(\rightarrow\) \textbf{build mode}:
        each output is read into a record, compositions are resolved, and a
        new database is written;
  \item \textbf{a mixture of both kinds} is a usage error (exit code 2).
\end{itemize}

The file extension is irrelevant: a database named \texttt{weird.dat} still
merges if its first non-blank line begins with \texttt{\{}, and a Gaussian
log named \texttt{weird.xyz} is still read as Gaussian output.

\subsection{Three families of flags}

Flags fall into three groups, and the distinction matters for what a run is
allowed to do.

\begin{description}
  \item[Manipulation flags] change which records survive, or add corrected
        values to them:
        \flg{--extract}, \flg{--nextract}, the cutoff flags, \flg{--uniq},
        \flg{--reacted}, \flg{--reacted2}, \flg{--sort}, \flg{--reverse},
        \flg{--select}, \flg{--glob}, \flg{--shuffle}, \flg{--seed},
        \flg{--noname}, \flg{--add-correction}, and \flg{--extrapolate}.
  \item[Print flags] ask for output columns or sections: \flg{--filename},
        \flg{--el}, \flg{--g}, \flg{--pop}, \flg{--formation},
        \flg{--hydration}, \flg{--movie}, and the rest.
  \item[Settings flags] configure derivations or presentation but do not
        count as a manipulation or a print: \flg{--temp}, \flg{--fc},
        \flg{--fscale}, \flg{--makereal}, \flg{--ignoreimag}, \flg{--unit},
        \flg{--conc}, \flg{--cnt}, \flg{--bavg}, \flg{--rh}, and
        \flg{--psolvent}.
\end{description}

A run must have a \emph{visible outcome}. With neither manipulation nor print
flags, \flg{--out} is required and the run behaves exactly like a collection
run. A run with a manipulation flag but neither \flg{--out} nor a print flag
is a usage error, because the manipulated records would be stored nowhere and
nothing would be printed.

\begin{note}
A settings flag on its own is not a visible outcome. For example,
\texttt{CJKQC db.jsonl} \flg{--temp} \texttt{350} is a usage error; add a print flag or
\flg{--out}.
\end{note}

\subsection{Overwriting, summaries, and exit codes}

A file that already exists is never replaced silently. Immediately before the
old file is replaced, \CJKQC{} asks:

\begin{verbatim}
db.jsonl already exists; overwrite? [y/N]
\end{verbatim}

Only \texttt{y} or \texttt{yes} (any casing, surrounding whitespace ignored)
proceeds. When standard input is not an interactive terminal, the question is
not asked and the replacement is refused with a message telling you to use
\flg{--fout}. \flg{--fout} forces the replacement without asking.

The completion line printed by the program tells you what happened:

\begin{verbatim}
collected 120 records into db.jsonl (2 failed, 1 truncated)
merged 40 records into merged.jsonl
filtered 17 records into kept.jsonl
wrote 17 frames into movie.xyz
\end{verbatim}

Exit codes are:

\begin{center}
\begin{tabular}{@{}cl@{}}
\toprule
Code & Meaning \\
\midrule
0 & success, nothing skipped \\
1 & fatal error (nothing written) or a merge that skipped corrupt lines \\
2 & usage error (bad flag, missing \flg{--out}, mixed inputs) \\
\bottomrule
\end{tabular}
\end{center}

Failed and truncated records from a recognized format do not change the exit
code: they are data, not failures of the run.

\subsection{Quick start}

Collect every Gaussian and ORCA output in the current directory into a new
database, then print the names, electronic energies, and Gibbs free energies
of the structures within 10 kcal/mol of the lowest electronic energy of their
composition group:

\begin{verbatim}
# 1. Build a database (the shell expands *.out)
CJKQC *.out -o all.jsonl

# 2. Print a table
CJKQC all.jsonl --filename --el --g

# 3. Filter, print, and store the survivors
CJKQC all.jsonl --cutr el 10 --filename --el --g -o near-minima.jsonl
\end{verbatim}

The next chapters explain each step in detail, beginning with how data enters
the database and how to read the resulting database file.

\section{Loading data}\label{sec:loading}

This chapter covers both directions of the data flow: how to turn
quantum-chemistry outputs into a database, and how to read a database
without programming. It also defines the standard units and the information
each record holds, because every later chapter refers to them.

\subsection{Collecting program outputs into a database}\label{sec:collecting}

Build mode is used when every positional input is a quantum-chemistry output
file. The command is simply:

\begin{verbatim}
CJKQC <output files>... -o <database.jsonl>
\end{verbatim}

For each input, in the order given, \CJKQC{}:

\begin{enumerate}
  \item opens the file and looks at its leading lines to decide how to read
        it;
  \item reads the head and tail windows of the file;
  \item extracts the main calculated values of the final job (or final
        complete xyz frame);
  \item resolves the composition from the record's name (the file stem);
  \item stores exactly one record.
\end{enumerate}

The output database is replaced only after the whole run succeeds: either the
new database is installed, or an error leaves any previous database intact. A
failed or truncated job from a recognized format is still written, with status
\fld{failed} or \fld{truncated} and a short explanation in its error entry.

\begin{warning}
An input that matches no supported format, or an empty input, is not a job. It
aborts the run with exit code 1 and nothing is written. If you want to mix
databases and program outputs, merge them in a separate invocation.
\end{warning}

\subsubsection*{Windows and very large files}

\CJKQC{} does not read the whole of a large output file. It reads a head
window (default 256 KiB) and a tail window (default 4 MiB), which is where
banners, route sections, final energies, geometries, and frequencies live.
Files smaller than the two windows together are read whole. If a record's
error says that a table was cut by a window seam, widen the windows:

\begin{verbatim}
CJKQC huge.out --head-bytes 1048576 --tail-bytes 8388608 -o db.jsonl
\end{verbatim}

These flags affect build-mode reads only; merging databases ignores them.
A value below 1 is a usage error.

\subsubsection*{Reading several files at once}

By default collection reads one file at a time. On a shared cluster
filesystem, or when many files are being read, you can ask \CJKQC{} to work
on several files at once:

\begin{verbatim}
CJKQC *.log *.out --parallel 8 -o db.jsonl
\end{verbatim}

Each file is still handled as a whole, and the resulting database is identical
to the one produced by a one-at-a-time run, regardless of the number you give.
A value below 1 is a usage error. Merging databases always reads one file at a
time and ignores \flg{--parallel} for reading; if the run also uses
\flg{--arbalign}, that filtering step still uses the setting
(Section~\ref{sec:arbalign}).

\subsubsection*{Composition from file names}

The record's \fld{name} is the file stem, for example \texttt{2am3sa-1conf}
from \texttt{2am3sa-1conf.log}. When the stem follows the JKQC naming
convention, collection also stores the composition. Section
\ref{sec:names} explains the grammar. A name that does not follow the
convention is kept exactly as written and simply carries no composition
information.

\subsection{Merging databases}\label{sec:merging}

When every positional input is a database, \CJKQC{} runs in merge mode:

\begin{verbatim}
CJKQC part1.jsonl part2.jsonl -o combined.jsonl
\end{verbatim}

Records are concatenated in argument order and copied exactly as they were
stored: there is no re-reading, no composition resolution, and no rewriting of
the records. Information written by a newer version of \CJKQC{} therefore
survives a merge through an older version. Corrupt lines are skipped, reported
by path and line number, and make the exit code 1, but the valid records are
still written. \flg{--out} may name one of the inputs; all inputs are read
before the output is replaced.

\subsection{Reading a database file as a non-programmer}\label{sec:jsonl}

A database is a plain text file. Each line holds one complete record, written
as a set of names and values inside curly braces. The rules are deliberately
simple:

\begin{itemize}
  \item one record per line;
  \item every complete line ends with a line break;
  \item blank lines are ignored when the file is read;
  \item the \texttt{.jsonl} ending is only a convention; the content decides
        whether a file is a database;
  \item a corrupt or partially written line is skipped and reported, so the
        rest of the file remains usable;
  \item records may repeat; duplicate removal is a filtering decision, not a
        storage decision.
\end{itemize}

In the real file, each record is written on one long line. For readability,
the same record is shown below spread over several lines. Information that is
not available is left out completely; you will not see the word
\texttt{null}.

\begin{verbatim}
{
  "status": "ok",
  "name": "1dma",
  "source_path": "/scratch/runs/xtb/1dma.log",
  "program": "xtb",
  "program_version": "6.7.1",
  "command": "xtb 1dma.xyz --gxtb --acc 0.01 -ohess",
  "electronic_energy": -38.51,
  "charge": 0,
  "multiplicity": 1,
  "geometry": {
    "symbols": ["N", "H"],
    "coordinates": [0.0, 0.0, 0.0, 0.0, 0.0, 1.01]
  },
  "vibrational_frequencies": [3421.1],
  "components": {"dma": 1}
}
\end{verbatim}

What to notice:

\begin{itemize}
  \item \fld{status} says whether the job finished. If this entry is missing,
        it is treated as \fld{ok}.
  \item \fld{name} is the identity used for matching corrections and monomer
        references. \fld{source_path} is the full path of the original output
        file, when it could be determined.
  \item \fld{geometry.symbols} and \fld{geometry.coordinates} form two
        matching lists: one element name per atom, and three coordinates per
        atom in the same order. The coordinates are listed atom by atom: $x,y,z$
        for the first atom, then $x,y,z$ for the second, and so on.
  \item \fld{components} lists each component code together with its count.
        The codes are written alphabetically, so the same composition always
        looks the same.
  \item A \fld{correction} entry appears only after a correction has been
        attached. It never replaces \fld{electronic_energy}.
\end{itemize}

The entries may appear in any order, and unknown entries are ignored. New
versions of \CJKQC{} only add optional entries; they never rename or
repurpose an existing one. An older database therefore remains readable, and
a newer one remains readable by older tools.

\subsection{Record entries and standard units}\label{sec:fields}

The following table lists every entry a record can hold. Energies,
coordinates, frequencies, dipoles, and polarizabilities are always stored in
the standard units shown; conversion happens only when values are printed.

\begin{longtable}{@{}L{5.4cm}L{2.2cm}L{5.9cm}@{}}
\toprule
\textbf{Entry} & \textbf{Kind of value} & \textbf{Meaning and standard unit} \\
\midrule
\endhead
\fld{status} & text & \fld{ok}, \fld{failed}, or \fld{truncated}. \\
\fld{error} & text & Plain-language reason for a failed or truncated job;
present whenever the status is \fld{failed}. \\
\fld{name} & text & Structure identity: the source file's stem, exactly as
written. \\
\fld{source_path} & text & Full path of the original output file. \\
\fld{program} & text & Program name: \fld{gaussian}, \fld{orca}, \fld{xtb},
\fld{mrcc}, or \fld{xyz}. \\
\fld{program_version} & text & Version string as printed by the program. \\
\fld{command} & text & The run's command line as echoed by the program. \\
\fld{electronic_energy} & number & Total electronic energy of the final
structure, in hartree. \\
\fld{scf_energy} & number & Converged SCF energy of the final structure, in
hartree. Used only by extrapolation. \\
\fld{charge} & whole number & Net charge, in elementary charges. \\
\fld{multiplicity} & whole number & Spin multiplicity, at least 1. \\
\fld{geometry} & two lists & \fld{symbols} (element names) and
\fld{coordinates} (angstrom, listed atom by atom). \\
\fld{vibrational_frequencies} & list & Genuine vibrational modes in
cm$^{-1}$; a negative value marks an imaginary mode. \\
\fld{gradient} & list & Cartesian gradient in hartree/\AA{}, listed atom by
atom, three values per atom. \\
\fld{dipole_moment} & list & Dipole vector $[x,y,z]$ in atomic units
($e\,a_0$). \\
\fld{polarizability} & list & The nine polarizability tensor values in
atomic units, in the order $xx, xy, xz, yx, yy, yz, zx, zy, zz$. The
isotropic value is calculated from them when printed. \\
\fld{mulliken_charges} & list & Per-atom Mulliken charges in elementary
charges, in geometry atom order. \\
\fld{rotational_symmetry_number} & whole number & Point-group rotational
symmetry number $\sigma$, at least 1. \\
\fld{components} & code and count & Resolved composition. \\
\fld{correction} & group of values & Attached single-point correction and its
origin; see Chapter~\ref{sec:corrections}. \\
\bottomrule
\end{longtable}

The standard units are summarized here once:

\begin{center}
\begin{tabular}{@{}ll@{}}
\toprule
Quantity & Unit stored in the database \\
\midrule
Energy, enthalpy, Gibbs, zero-point energy & hartree \\
Coordinates and distances & \AA{} (angstrom) \\
Gradient & hartree/\AA{} \\
Vibrational frequencies & cm$^{-1}$ \\
Dipole moment & atomic units, $e\,a_0$ \\
Polarizability & atomic units, $e^2 a_0^2/E_\mathrm{h}$ \\
Charge and Mulliken charges & elementary charge $e$ \\
Mass & unified atomic mass unit, u \\
Multiplicity and rotational symmetry number & whole numbers \\
\bottomrule
\end{tabular}
\end{center}

Printing in kcal/mol is a presentation choice made with \flg{--unit}; it
never changes the database.

\subsection{Names and compositions}\label{sec:names}

The JKQC naming convention encodes a composition in the file stem:

\begin{center}
\begin{tabular}{@{}ll@{}}
\toprule
Name & Resolved components \\
\midrule
\texttt{1dma} & \texttt{\{"dma": 1\}} \\
\texttt{2cha} & \texttt{\{"cha": 2\}} \\
\texttt{1pinic1sa} & \texttt{\{"pinic": 1, "sa": 1\}} \\
\texttt{2am3sa-1conf} & \texttt{\{"am": 2, "sa": 3\}} \\
\texttt{1am1p} & \texttt{\{"am": 1, "p": 1\}} \\
\texttt{h2o} & no composition (does not match the grammar) \\
\texttt{movie} & no composition \\
\bottomrule
\end{tabular}
\end{center}

The grammar is:

\begin{verbatim}
name        := composition ( "-" label )?
composition := ( count code )+        (at least one pair)
count       := digit+                  (1 to 4294967295)
code        := letter+                 (plain letters, either case)
label       := any text                (never read, never stored)
\end{verbatim}

Details worth remembering:

\begin{itemize}
  \item The name splits at its \emph{first} hyphen. Everything after it is a
        free-form label and is ignored.
  \item Counts allow leading zeros: \texttt{02am} means two \texttt{am}.
  \item Codes are case-sensitive: \texttt{am} and \texttt{Am} are different
        components.
  \item Repeated codes sum: \texttt{1am1am} resolves to
        \texttt{\{"am": 2\}}.
  \item A count of zero, a special character, a digit inside a code, or an
        implicit count such as \texttt{am2sa} makes the name unresolvable.
        This is never an error; the record simply has no \fld{components}.
\end{itemize}

Later features use the stored composition, not the raw name. This is
why a database written once can be filtered and post-processed without
re-reading file names.

\subsection{Record statuses}\label{sec:statuses}

\begin{description}
  \item[\fld{ok}] The job converged and its values are final.
  \item[\fld{failed}] The job did not finish. Some values may still be
        present, and \fld{error} explains what happened.
  \item[\fld{truncated}] The output file ends in the middle of the run. Some
        values may still be present.
\end{description}

The filtering steps drop all non-\fld{ok} records before anything else and
report their count as a warning. A print-only run prints the
records as loaded, so a failed record can still appear in a table unless you
filter it out.

\section{Printing}\label{sec:printing}

Printing works one column at a time: every printable quantity has its own
flag, and there is no default selection. If you pass no print flag, no table
is printed.
Columns always appear in the same fixed order, no matter what order the flags
were given on the command line.

\begin{verbatim}
CJKQC db.jsonl --filename --el --g --rg
\end{verbatim}

\subsection{How the table is formatted}\label{sec:table-format}

The table is plain text output: a header row of column labels,
then one row per record in the order the run produces. Columns are as wide as
their widest cell, with one space between columns and no trailing whitespace.
Text columns are left-aligned; numeric columns are right-aligned; header
cells align with their columns.

An undefined or absent value prints as a single dash \texttt{-}. A
list-valued column with no values also prints \texttt{-}. Numbers keep fixed
per-column formats, so table columns line up and can be pasted into other
tools.

The fixed column order is:

\begin{center}
\small
\begin{tabular}{@{}llll@{}}
\toprule
\multicolumn{4}{c}{\textbf{Fixed print order}} \\
\midrule
\fld{filename} & \fld{name} & \fld{path} & \fld{status} \\
\fld{program} & \fld{command} & \fld{correction} & \fld{el} \\
\fld{zpe} & \fld{h} & \fld{s} & \fld{g} \\
\fld{elcorr} & \fld{hcorr} & \fld{gcorr} & \fld{pop} \\
\fld{popel} & \fld{popelcorr} & \fld{popgcorr} & \fld{rg} \\
\fld{mass} & \fld{dip} & \fld{dips} & \fld{pol} \\
\fld{pols} & \fld{charge} & \fld{multiplicity} & \fld{rsn} \\
\fld{natoms} & \fld{lf} & \fld{freqs} & \fld{nfreq} \\
\fld{nimag} & \fld{rots} & \fld{components} & \\
\bottomrule
\end{tabular}
\end{center}

\subsection{Identity and source columns}\label{sec:identity-columns}

\begin{longtable}{@{}L{3.1cm}L{2.4cm}L{7.7cm}@{}}
\toprule
\textbf{Flag} & \textbf{Header} & \textbf{What it prints} \\
\midrule
\endhead
\flg{--filename} & \fld{filename} & The record's \fld{name}: the stem of the
source file name. \\
\flg{--name} & \fld{name} & The composition written as alphabetical pairs
without separators, for example \texttt{2am3sa}. Undefined without
components. \\
\flg{--path} & \fld{path} & The record's \fld{source_path}, exactly as
stored. \\
\flg{--status} & \fld{status} & \fld{ok}, \fld{failed}, or \fld{truncated}. \\
\flg{--program} & \fld{program} & The program name and version joined by one
space, for example \texttt{xtb 6.7.1}; the name alone when the version is
absent. \\
\flg{--command} & \fld{command} & The stored command line in double quotes,
because it contains spaces. An absent command prints a bare \texttt{-}. \\
\flg{--correction} & \fld{correction} & Correction source details: the kind, then
the program and version, then each command in double quotes. \\
\flg{--components} & \fld{components} & The resolved composition as pairs with
plus signs, for example \texttt{2am+3sa}. \\
\bottomrule
\end{longtable}

\subsection{Energy columns}\label{sec:energy-columns}

The energy columns share one format: hartree with 10 decimals by default, or
kcal/mol with 4 decimals under \flg{--unit}. The list includes derived
thermochemistry and corrected values; each is explained in
Chapter~\ref{sec:thermo} and Chapter~\ref{sec:corrections}.

\begin{center}
\begin{tabular}{@{}lll@{}}
\toprule
\textbf{Flag} & \textbf{Header} & \textbf{Quantity} \\
\midrule
\flg{--el} & \fld{el} & Electronic energy \\
\flg{--zpe} & \fld{zpe} & Harmonic zero-point energy \\
\flg{--h} & \fld{h} & Enthalpy at \flg{--temp} \\
\flg{--s} & \fld{s} & Entropy term $T\,S$ at \flg{--temp} \\
\flg{--g} & \fld{g} & Gibbs free energy at \flg{--temp} \\
\flg{--elcorr} & \fld{elcorr} & Attached correction energy \\
\flg{--hcorr} & \fld{hcorr} & Enthalpy from the correction energy \\
\flg{--gcorr} & \fld{gcorr} & Gibbs free energy from the correction energy \\
\bottomrule
\end{tabular}
\end{center}

\flg{--unit} is a simple on/off switch: adding it alone converts every
hartree-valued print quantity to kcal/mol, every dipole value to debye,
and every polarizability value to \AA$^3$. It takes no value, and a value
after it is a usage error. Without it, energies print in hartree and
dipoles and polarizabilities in atomic units. Relative
cutoffs (Chapter~\ref{sec:filtering}) already compare in kcal/mol and are not
affected by \flg{--unit}, and the dipole cutoff remains in atomic units.

\subsection{Structure and property columns}\label{sec:property-columns}

\begin{longtable}{@{}L{2.9cm}L{2.5cm}L{7.8cm}@{}}
\toprule
\textbf{Flag} & \textbf{Header} & \textbf{What it prints} \\
\midrule
\endhead
\flg{--rg} & \fld{rg} & Mass-weighted radius of gyration in \AA{}, 3
decimals. \\
\flg{--mass} & \fld{mass} & Total mass in u, 3 decimals. \\
\flg{--dip} & \fld{dip} & Dipole magnitude in atomic units, or debye under
\flg{--unit}, 3 decimals. \\
\flg{--dips} & \fld{dips} & The dipole vector's $x,y,z$ components in atomic
units, or debye under \flg{--unit}, 3 decimals each, comma-separated
without spaces. \\
\flg{--pol} & \fld{pol} & Isotropic polarizability in atomic units, or
\AA$^3$ under \flg{--unit}, 3 decimals. \\
\flg{--pols} & \fld{pols} & The nine polarizability tensor values in atomic
units, or \AA$^3$ under \flg{--unit}, 3 decimals each, comma-separated
without spaces, in the order $xx, xy, xz, yx, yy, yz, zx, zy, zz$. \\
\flg{--charge} & \fld{charge} & Net charge as a whole number. \\
\flg{--multiplicity} & \fld{multiplicity} & Spin multiplicity as a whole
number. \\
\flg{--rsn} & \fld{rsn} & Rotational symmetry number as a whole number. \\
\flg{--natoms} & \fld{natoms} & Geometry atom count. \\
\bottomrule
\end{longtable}

\subsection{Frequency columns}\label{sec:frequency-columns}

\begin{longtable}{@{}L{2.9cm}L{2.5cm}L{7.8cm}@{}}
\toprule
\textbf{Flag} & \textbf{Header} & \textbf{What it prints} \\
\midrule
\endhead
\flg{--lf} & \fld{lf} & The lowest stored vibrational frequency in
cm$^{-1}$, 1 decimal. Raw values, so an imaginary mode prints negative. \\
\flg{--freqs} & \fld{freqs} & Every stored vibrational frequency in stored
order, 1 decimal each, comma-separated without spaces. Raw signs are kept. \\
\flg{--nfreq} & \fld{nfreq} & Number of stored vibrational frequencies. \\
\flg{--nimag} & \fld{nimag} & Number of negative (imaginary) frequencies. \\
\flg{--rots} & \fld{rots} & Rotational constants in cm$^{-1}$, descending,
6 decimals each, comma-separated without spaces. \\
\bottomrule
\end{longtable}

\begin{note}
The frequency columns show the \emph{raw stored values}. They are not
changed by \flg{--makereal}, \flg{--ignoreimag}, or \flg{--fscale}; those
settings affect only the thermochemistry derivation
(Section~\ref{sec:freq-manipulation}). A saddle point therefore still prints
a negative \fld{lf} and \fld{nimag} even when its Gibbs energy is made
defined.
\end{note}

\subsection{Boltzmann population columns}\label{sec:populations}

The \flg{--pop} flag prints the fraction of the thermal ensemble that each
printed structure occupies. It is repeatable and takes an optional level:

\begin{center}
\begin{tabular}{@{}lll@{}}
\toprule
\textbf{Invocation} & \textbf{Column} & \textbf{Energy basis} \\
\midrule
\flg{--pop} or \flg{--pop g} & \fld{pop} & Gibbs free energy at \flg{--temp} \\
\flg{--pop el} & \fld{popel} & Electronic energy \\
\flg{--pop elcorr} & \fld{popelcorr} & Attached correction energy \\
\flg{--pop gcorr} & \fld{popgcorr} & Gibbs free energy from the correction \\
\bottomrule
\end{tabular}
\end{center}

Populations are computed over exactly the records the run prints, grouped by
composition, or into one global group under \flg{--noname}. For a group and a
level with values $v_i$ in hartree,

\begin{equation}
p_i = \frac{\exp\!\left(-\dfrac{v_i - v_{\min}}{k_B T}\right)}
           {\sum_j \exp\!\left(-\dfrac{v_j - v_{\min}}{k_B T}\right)},
\end{equation}

where $v_{\min}$ is the group's minimum defined value and $T$ is
\flg{--temp}. Only status \fld{ok} records with a defined level value
contribute. A record with an undefined value prints \texttt{-}; if no member
of the group defines the level, every member prints \texttt{-}. Populations
are dimensionless probabilities printed with 4 decimals, and \flg{--unit}
never affects them.

\begin{warning}
A space-separated value after \flg{--pop} is consumed by the flag. That means
\texttt{CJKQC} \flg{--pop} \texttt{db.jsonl} tries to use \texttt{db.jsonl} as a population
level and fails. Put \flg{--pop} after the positional inputs, or write
\flg{--pop=g}.
\end{warning}

\subsection{Sections and file output}\label{sec:sections-output}

Several print flags produce a section after the main table instead of a
column:

\begin{itemize}
  \item \flg{--formation} prints the formation-energy section
        (Chapter~\ref{sec:formation});
  \item \flg{--formation-corr} prints the corrected formation section;
  \item \flg{--hydration} and \flg{--hydration-corr} print water
        distributions (Chapter~\ref{sec:solvation});
  \item \flg{--solvation CODE} and \flg{--solvation-corr CODE} print the
        distribution for another solvent.
\end{itemize}

The standard-output order is always: main table, raw formation section,
corrected formation section, hydration/solvation sections (raw before
corrected), then the completion summary. A run with only section flags prints
the section and no main table.

\flg{--movie} writes the final structures into \texttt{movie.xyz}, one xyz
frame per \fld{ok} record that carries a geometry, in the order the table
prints. It is a print flag and a visible outcome on its own. Records without
a geometry, or with a non-\fld{ok} status, are skipped and counted in a
standard-error note. The file is always named \texttt{movie.xyz} in the
current working directory and is subject to the same overwrite confirmation
as \flg{--out}.

\flg{--write-xyz} writes one separate xyz file for each \fld{ok} record that
has a geometry and a usable \fld{filename}. The file is named after the
record: a record whose \fld{filename} is \texttt{1sa-3xstr} becomes
\texttt{1sa-3xstr.xyz} in the current working directory. Each file has the
same one-structure content as the corresponding \flg{--movie} frame. It is a
print flag and a visible outcome on its own, and it can be combined with
\flg{--out} and \flg{--movie}.

Records that are not \fld{ok}, have no geometry, have no name, or have a name
that cannot be used as a file name are skipped and counted in a
standard-error note. Names that would produce the same file, including names
that differ only by upper and lower case, stop the run as a usage error
before anything is written. A record's file also must not have the same name
as \texttt{movie.xyz} in a \flg{--movie} run, or as the \flg{--out} database
in a run that writes one. The existing files are covered by one overwrite
confirmation, and \flg{--fout} forces their replacement.

When \flg{--out} is also used, the database is written first, then the
separate xyz files, then \texttt{movie.xyz}. The completion line is
\texttt{wrote <n> xyz files into the current directory}.

\subsection{A complete example}\label{sec:printing-example}

\begin{verbatim}
CJKQC db.jsonl --filename --program --el --g --lf --nimag --pop
\end{verbatim}

This prints the file stem, program name and version, electronic energy,
Gibbs free energy, lowest frequency, imaginary-mode count, and Gibbs-based
population. The exact widths depend on the data; the layout is illustrated
below with hypothetical values.

\begin{verbatim}
filename  program     el            g             lf     nimag  pop
1dma      xtb 6.7.1   -38.5100000000 -38.4987654321 3421.1 0      0.7321
1dma-2    xtb 6.7.1   -38.5090000000 -38.4970000000 3402.7 0      0.2679
\end{verbatim}

Everything printed here is derived from the database at the run's settings.
Nothing in the table was stored when the database was built.

\section{Filtering, sorting, and duplicate removal}\label{sec:filtering}

Filtering turns a loaded set of records into a smaller, ordered set. Most steps
only read the records; renaming component codes rewrites a record's
composition. No step changes the input database on disk. A filtering run writes
the records that remain with \flg{--out}, prints them with print flags, or both.
Print flags alone do not filter: a print-only run shows every loaded record in
file order.

\subsection{Order of the filtering steps}\label{sec:pipeline-order}

The steps always run in this fixed order, each within the current composition
groups:

\begin{enumerate}
  \item drop records whose status is not \fld{ok};
  \item rename component codes with \flg{--rename}, in the order the flags were
        given;
  \item \flg{--extract} or \flg{--nextract};
  \item cutoffs, in the order the cutoff flags were given;
  \item \flg{--uniq};
  \item \flg{--reacted} or \flg{--reacted2};
  \item \flg{--sort}, then \flg{--reverse};
  \item \flg{--select};
  \item \flg{--arbalign};
  \item \flg{--glob};
  \item \flg{--shuffle}.
\end{enumerate}

The non-\fld{ok} count is reported as a warning:
\texttt{excluded <k> records with status failed or truncated}.

When \flg{--uniq} is used with a list of properties, or when \flg{--select}
or \flg{--arbalign} is used, and no \flg{--sort} is given, the records are
first sorted by \fld{el} ascending. \flg{--uniq dup} does not trigger this
automatic sort.

\subsection{Grouping}\label{sec:grouping}

Most steps work per \emph{composition group}: records with the same
composition form a group, and records without composition information form
one separate unnamed group. Groups are processed in first-appearance order;
within a step, the overall record order is preserved.

\flg{--noname} collapses everything into one global group. Under
\flg{--noname}, \flg{--extract} and \flg{--nextract} compare their patterns
directly against each record's \fld{name} (the same text the \fld{filename}
column prints), with no composition grammar, ranges, wildcards, or zero-count
rules. \flg{--arbalign} is the exception: it always compares records only
when they have the same stored composition, so \flg{--noname} does not make
different compositions comparable.

\subsection{Renaming component codes}\label{sec:rename}

\flg{--rename OLD NEW} moves the count of one component code to another code.
Both codes are plain letters, and upper and lower case are different codes. The
flag may be repeated; renames apply in the order given and chain together, so
\flg{--rename p a --rename a sa} first moves \texttt{p} counts to \texttt{a}
and then moves those \texttt{a} counts on to \texttt{sa}.

The rename applies to every record that has a stored composition:

\begin{itemize}
  \item If the record does not contain \texttt{OLD}, its composition is left
        unchanged.
  \item If \texttt{OLD} and \texttt{NEW} are the same code, nothing changes.
  \item Otherwise the count of \texttt{OLD} is added to the count of
        \texttt{NEW}. For example, a record with 3 \texttt{sa}, 1 \texttt{am},
        and 1 \texttt{p} becomes 4 \texttt{sa} and 1 \texttt{am} under
        \flg{--rename p sa}.
  \item A record with no stored composition is left unchanged.
\end{itemize}

Renaming changes the stored composition, not the record's file name or source
path. The \fld{filename} column still prints the original file stem, while the
\fld{name} and \fld{components} columns and all composition groups show the new
composition. Because renaming is the first filtering step after failed records
are dropped, later steps such as \flg{--extract} and relative cutoffs work with
the new codes.

\subsection{Extracting compositions}\label{sec:extract}

\flg{--extract QUERY} keeps records whose composition matches the query;
\flg{--nextract QUERY} keeps the records that do not. A query is a
comma-separated list of patterns; a record matches if any pattern matches.

A pattern is built from these pieces:

\begin{verbatim}
query   := pattern ( "," pattern )*
pattern := ( pair | range | choice | bare-star | count-star
             | count-star-code | star-code )+
pair    := count code
range   := count "-" count code          (inclusive count expansion)
choice  := "(" count ( "," count )* ")" code
bare-star       := "*"
count-star      := count "*"
count-star-code := count "*" code
star-code       := "*" code
count   := digit+ (decimal, at most 4294967295)
code    := letter+ (plain letters, case matters)
\end{verbatim}

A pair \texttt{(count, code)} with a positive count requires that exact count
for the code. A count of zero requires the code to be absent. Codes the
pattern does not mention are \emph{extras}, governed by the wildcard:

\begin{center}
\begin{tabular}{@{}ll@{}}
\toprule
\textbf{Form} & \textbf{Extra records allowed} \\
\midrule
no wildcard & none: the pattern is an exact composition match \\
\texttt{*} & any composition \\
\texttt{N*} & extras whose counts sum to exactly $N$ \\
\texttt{*code} & only the named codes, at any counts \\
\texttt{N*code} & only the named codes, with counts summing to exactly $N$ \\
\bottomrule
\end{tabular}
\end{center}

Examples:

\begin{center}
\begin{tabular}{@{}ll@{}}
\toprule
\textbf{Query} & \textbf{Matches} \\
\midrule
\texttt{2am3sa} & exactly 2 \texttt{am} and 3 \texttt{sa}, no extras \\
\texttt{2am3sa,1am1sa} & either of the two compositions \\
\texttt{1am0-2w} & \texttt{1am}, \texttt{1am1w}, \texttt{1am2w} \\
\texttt{(2,3)am1w} & \texttt{2am1w} or \texttt{3am1w} \\
\texttt{2am*w} & 2 \texttt{am} plus any number of \texttt{w} \\
\texttt{2am2*w} & 2 \texttt{am} plus \texttt{w} counts summing to 2 \\
\bottomrule
\end{tabular}
\end{center}

Further rules:

\begin{itemize}
  \item A range expands inclusively. A range with more than 1000 count values,
        or a pattern expanding beyond 10000 patterns, is a usage error.
  \item A range whose low end exceeds its high end is a usage error.
  \item Anything after a \texttt{-} that a range did not consume is treated
        as a label comment and ignored, so a full JKQC file stem can be
        pasted directly.
  \item An underscore is not a valid pattern character; a pattern containing
        one is a usage error.
  \item Records without \fld{components} behave as if their composition were
        empty: they match only patterns with no pairs, such as \texttt{*}.
\end{itemize}

\subsection{Cutoffs}\label{sec:cutoffs}

Each cutoff compares a record's property value with a number. The ten flags
are:

\begin{center}
\begin{tabular}{@{}lll@{}}
\toprule
\textbf{Absolute flag} & \textbf{Comparison} & \textbf{Mode} \\
\midrule
\flg{--cut} & $\le$ & absolute \\
\flg{--filter-lt} & $<$ & absolute \\
\flg{--pass} & $>$ & absolute \\
\flg{--filter-ge} & $\ge$ & absolute \\
\flg{--is} & $=$ & absolute \\
\flg{--isnot} & $\ne$ & absolute \\
\midrule
\flg{--cutr} & $\le$ & relative \\
\flg{--rel-filter-lt} & $<$ & relative \\
\flg{--passr} & $>$ & relative \\
\flg{--rel-filter-ge} & $\ge$ & relative \\
\bottomrule
\end{tabular}
\end{center}

They are written \flg{--cut PROPERTY VALUE} and so on, and they may repeat;
cutoffs chain in command order.

\begin{center}
\begin{tabular}{@{}lll@{}}
\toprule
\textbf{Property} & \textbf{Value} & \textbf{Units} \\
\midrule
\fld{el} & electronic energy & hartree absolute; kcal/mol relative \\
\fld{g} & Gibbs free energy & hartree absolute; kcal/mol relative \\
\fld{elcorr} & attached correction energy & hartree absolute; kcal/mol relative \\
\fld{gcorr} & Gibbs from the correction & hartree absolute; kcal/mol relative \\
\fld{rg} & radius of gyration & \AA{} \\
\fld{lf} & lowest stored frequency & cm$^{-1}$ \\
\fld{dip} & dipole magnitude & atomic units \\
\fld{mass} & total mass & u \\
\fld{charge} & net charge & counts \\
\fld{multiplicity} & spin multiplicity & counts \\
\fld{natoms} & geometry atom count & counts \\
\fld{nfreq} & stored frequency count & counts \\
\fld{nimag} & negative frequency count & counts \\
\bottomrule
\end{tabular}
\end{center}

In relative mode the group minimum of the property is subtracted first (over
defined values), then the comparison is applied. Thus
\flg{--cutr el 10} keeps records within 10 kcal/mol of their group's lowest
electronic energy. Under \flg{--noname} the minimum is global.

A record whose property is undefined fails $<$, $\le$, $>$, $\ge$, and $=$,
and passes $\ne$. This mirrors the original tool's treatment of missing
values.

\begin{warning}
Absolute energy cutoffs are in hartree and relative energy cutoffs in
kcal/mol, regardless of whether \flg{--unit} is used for printing. A value
such as \flg{--cut el -38.5} is hartree; \flg{--cutr el 10} is kcal/mol.
\end{warning}

\subsection{Sorting, selecting, and shuffling}\label{sec:sorting}

\flg{--sort PROPERTY} sorts ascending; \flg{--reverse} reverses the
orientation. The possible properties are:

\begin{center}
\begin{tabular}{@{}ll@{}}
\toprule
\textbf{Property} & \textbf{Orders by} \\
\midrule
\fld{el} & electronic energy \\
\fld{g} & Gibbs free energy at the run's conditions \\
\fld{elcorr} & attached correction energy \\
\fld{gcorr} & Gibbs free energy from the correction \\
\fld{rg} & radius of gyration \\
\fld{mass} & total mass \\
\fld{filename} & the record's \fld{name}; nameless records first \\
\bottomrule
\end{tabular}
\end{center}

A record whose numeric sorting value is undefined sorts last ascending and
first descending. \flg{--reverse} alone, with no sort in effect, changes
nothing but still counts as a manipulation.

\flg{--select N} keeps the first $N$ records of each group after sorting.
Groups shorter than $N$ pass whole. \flg{--glob} keeps, per composition
group, the record with the lowest Gibbs free energy at the run's conditions;
ties keep the first in order. Records with undefined Gibbs are dropped.
Under \flg{--noname}, \flg{--glob} keeps the one global minimum.

\flg{--shuffle} puts the final result in a random order. The order is driven
by a fixed starting value, \flg{--seed} (default 0), so the same seed always
gives the same shuffle.

\subsection{Uniqueness}\label{sec:uniqueness}

There are two ways to remove duplicate structures: by name or by recorded
property values, using \flg{--uniq}, and by three-dimensional geometry, using
\flg{--arbalign} (Section~\ref{sec:arbalign}). \flg{--uniq} removes
redundant structures within each composition group. Both modes are stable:
the first occurrence in the current order survives.

\begin{description}
  \item[\flg{--uniq dup}] Drops records that share a \fld{name} with an earlier
        record. Records without a name are never dropped as duplicates.
  \item[\flg{--uniq PROPERTIES}] Drops records whose rounded set of property
        values matches an earlier record in the same group.
\end{description}

The property list is a comma-separated list of \fld{metric[decimals]} items,
for example \flg{--uniq el2.4,rg}. The decimal count is optional; defaults are
\fld{el} 3, \fld{g} 3, \fld{rg} 2, \fld{mass} 2, and \fld{dip} 1. A fractional
decimal count is accepted. Two values agree at $d$ decimals exactly when

\begin{equation}
\lfloor a \cdot 10^d \rfloor = \lfloor b \cdot 10^d \rfloor,
\end{equation}

with floor toward $-\infty$. Undefined metrics are treated as equal to each
other, so two records that both lack the metric collapse.

Because the automatic sort is by \fld{el} ascending, writing
\flg{--uniq el2.4} without \flg{--sort} keeps the lowest-electronic-energy
member of each group of duplicates. Add \flg{--sort g} first if you want the
lowest-Gibbs representative instead.

\subsection{Removing duplicate geometries}\label{sec:arbalign}

Two conformers can have different names and slightly different energies but
still be the same three-dimensional structure. \flg{--arbalign THRESHOLD}
removes such duplicates by comparing geometries. The threshold is a distance
in \AA{}: a record is dropped when its best-fit root-mean-square distance
(RMSD) to an earlier surviving record of the same composition is strictly
less than the threshold. The RMSD is a size in \AA{} that measures how far
apart two structures are after the best possible rotation and atom matching.

Every structure is first turned so that its principal axes (the natural axes
of its mass distribution) line up. This always happens; there is no separate
flag to turn it on. The comparison then tries different atom orderings,
rotations, and mirror-image arrangements, so the order in which atoms appear
in the file and the molecule's orientation do not matter.

The step runs after \flg{--sort} and \flg{--select} and before
\flg{--glob} and \flg{--shuffle}. When no \flg{--sort} is given, \CJKQC{}
first sorts by \fld{el} ascending, so the lowest-electronic-energy member of
a duplicate set survives. Add \flg{--sort g} first if you want the
lowest-Gibbs member instead.

Further rules:

\begin{itemize}
  \item The comparison is strict: a pair exactly at the threshold is kept.
        A threshold of 0 therefore removes nothing.
  \item A negative, non-numeric, or infinite threshold is a usage error.
  \item Records without a usable geometry are kept and reported on the error
        output as
        \texttt{arbalign: kept <k> records without usable geometry}.
  \item A pair whose structures have different numbers or kinds of atoms is
        not compared.
  \item The comparison is done two records at a time: once a record is
        removed, it cannot remove another record. The first surviving record
        in the current order wins.
  \item \flg{--uniq} and \flg{--arbalign} can be combined. \flg{--uniq}
        runs before the sort; then the sort, \flg{--select}, and
        \flg{--arbalign} follow.
  \item \flg{--parallel} also controls how many geometry comparisons run at
        once. The survivors are the same at every setting.
\end{itemize}

\begin{note}
If you use \flg{--arbalign}, please cite the ArbAlign paper:
B.~Temelso, J.~M. Mabey, T.~Kubota, N.~Appiah-Padi, and G.~C. Shields,
``ArbAlign: A Tool for Optimal Alignment of Arbitrarily Ordered Isomers
Using the Kuhn--Munkres Algorithm,'' \emph{J. Chem. Inf. Model.}
\textbf{57} (2017), 1045--1054,
\href{https://doi.org/10.1021/acs.jcim.6b00546}{doi:10.1021/acs.jcim.6b00546}.
\end{note}

For example, keep the lowest-energy member of each group of structures that
are within 0.5 \AA{} RMSD of each other:

\begin{verbatim}
CJKQC db.jsonl --sort el --arbalign 0.5 \
  --filename --el --out unique.jsonl
\end{verbatim}

\subsection{Reacted and exploded structures}\label{sec:reacted}

Within a composition group, most conformers share a fragment fingerprint. A
structure whose optimization reacted or exploded has a different fingerprint.
\CJKQC{} builds each fingerprint from the geometry and uses the group's most
frequent fingerprint as the consensus.

Two atoms are bonded when their distance is below a pair-specific cutoff:

\begin{center}
\begin{tabular}{@{}lc@{}}
\toprule
\textbf{Pair} & \textbf{Cutoff (\AA{})} \\
\midrule
C--N (either order) & 1.75 \\
N--N & 1.50 \\
S--O or O--S & 1.90 \\
O--N or N--O & 1.90 \\
H--H & 0.80 \\
any other pair involving H & never bonded \\
every other heavy--heavy pair & 2.00 \\
\bottomrule
\end{tabular}
\end{center}

Each connected component's fragment formula is its element symbols sorted
alphabetically and concatenated. A record's signature is the sorted list of
its fragment formulas joined with \texttt{+}. For example, a bonded C--N
fragment plus a lone H is \texttt{CN+H}. A record without a geometry has the
signature \texttt{(none)} and participates like any other signature.

\flg{--reacted} keeps the records whose signature equals the group's
consensus; \flg{--reacted2} keeps exactly the deviants. Ties are resolved in a
fixed order: highest count first, then earliest record order.

\subsection{Writing the survivors}\label{sec:filter-out}

A manipulation run writes its survivors with \flg{--out} and reports:

\begin{verbatim}
filtered 17 records into kept.jsonl
\end{verbatim}

Correction and extrapolation add corrected energies to the records before
filtering starts, so cutoffs, sorts, and sections see those values
(Chapter~\ref{sec:corrections}). The overwrite confirmation and \flg{--fout}
rules are the same as in collection.

\subsection{A worked filtering example}\label{sec:filtering-example}

Keep only one conformer per rounded Gibbs/radius signature for each
composition, discard structures more than 5 kcal/mol above the group minimum,
sort by Gibbs energy, and print the lowest 3 per group:

\begin{verbatim}
CJKQC db.jsonl \
  --cutr g 5 \
  --uniq g2,rg1 \
  --sort g \
  --select 3 \
  --filename --name --g --rg --pop
\end{verbatim}

The steps apply in the fixed order: the relative Gibbs cutoff, then
uniqueness, then the sort, then the per-group limit. Populations are computed
over the final printed set, so they renormalize over the three survivors of
each group.

\section{Thermodynamics}\label{sec:thermo}

All thermochemistry is derived on demand. A record stores the electronic
energy, geometry, frequencies, multiplicity, and rotational symmetry number;
zero-point energy, enthalpy, entropy, and Gibbs free energy are recomputed
from those stored values under the settings of the current run. Program-printed
thermochemistry in the output file is never read or stored. This means a
database never needs to be rebuilt when you change the temperature, the
frequency scale factor, the low-frequency treatment, or the energy unit.

The default model is the ideal-gas rigid-rotor model with Grimme's
quasi-RRHO vibrational entropy, following the damping scheme used by ORCA,
xtb, and GoodVibes: below a reference frequency, the harmonic-oscillator
entropy is interpolated toward a free-rotor expression. ORCA~6 is used as the
reference for the implementation, and the derived values reproduce ORCA's
printed quasi-RRHO thermochemistry (Section~\ref{sec:thermo-reference}). The
reference frequency defaults to 100 cm$^{-1}$ and is changed with
\flg{--fc}.

\subsection{What is needed for a defined result}\label{sec:definedness}

A record has defined thermochemistry exactly when all of the following hold:

\begin{itemize}
  \item \fld{electronic_energy} is present;
  \item \fld{geometry} is present and non-degenerate (the middle principal
        moment of inertia is positive);
  \item every geometry symbol is in the atomic-weight table;
  \item no raw frequency is negative, unless \flg{--makereal} or
        \flg{--ignoreimag} handles imaginary modes;
  \item the species is monatomic, or it has at least one participating
        vibrational mode;
  \item every derived term is finite.
\end{itemize}

A non-monatomic record without frequencies, such as a single-point output or
a plain xyz file, has no thermochemistry. An undefined derivation prints
\texttt{-}, fails every cutoff comparison except $\ne$, sorts after defined
values, and is treated as an undefined metric by uniqueness. This is normal
and not an error.

\subsection{Conditions}\label{sec:thermo-conditions}

\begin{center}
\begin{tabularx}{\textwidth}{@{}L{2.7cm}L{2.0cm}L{2.7cm}X@{}}
\toprule
\textbf{Flag} & \textbf{Default} & \textbf{Constraint} & \textbf{Effect} \\
\midrule
\flg{--temp} & 298.15 K & $>0$ & Temperature $T$ \\
\flg{--fc} & 100 cm$^{-1}$ & $\ge 0$ & Quasi-harmonic reference frequency;
0 selects pure harmonic entropy \\
\flg{--fscale} & 1.0 & $>0$ & Multiplies every vibrational frequency before
the vibrational formulas \\
\flg{--makereal} & off & --- & Treats an imaginary mode as its absolute
value \\
\flg{--ignoreimag} & off & conflicts with \flg{--makereal} & Drops imaginary
modes from the vibrational sums \\
\bottomrule
\end{tabularx}
\end{center}

These are settings flags. They configure derivations, never count as
manipulations or print flags, and a run that passes only them still needs
\flg{--out} or a print flag.

\subsection{Frequency manipulation}\label{sec:freq-manipulation}

\subsubsection*{Participation}

After imaginary-mode handling, a mode participates in the vibrational sums
only if its frequency is greater than 1.0 cm$^{-1}$. This removes the
near-zero translational and rotational contaminant modes that some programs
print. A stored \texttt{-0.0} never participates. A genuine soft mode above
1.0 cm$^{-1}$ does participate, even if it is very low.

\subsubsection*{Scaling}

For every participating mode, the derivation uses

\begin{equation}
\tilde{\nu}'_i = \texttt{fscale} \cdot \tilde{\nu}_i .
\end{equation}

All vibrational formulas use $\tilde{\nu}'_i$. The 1.0 cm$^{-1}$
participation test is applied to the unscaled frequency, after imaginary-mode
handling.

\subsubsection*{Imaginary modes}

A negative raw frequency marks a saddle-point mode. The three behaviors are:

\begin{center}
\begin{tabular}{@{}L{3.2cm}L{4.1cm}L{6.0cm}@{}}
\toprule
\textbf{Setting} & \textbf{Treatment} & \textbf{Result} \\
\midrule
neither flag & no imaginary mode allowed & any negative frequency makes the
whole derivation undefined; all thermochemistry columns print \texttt{-} \\
\flg{--makereal} & each negative frequency is replaced by its absolute value
& the mode then follows the ordinary participation test and formulas \\
\flg{--ignoreimag} & each negative frequency is dropped before the
participation test & the mode contributes to no vibrational term; a
non-monatomic record with no remaining mode is undefined \\
\bottomrule
\end{tabular}
\end{center}

Neither flag edits the stored record. The \fld{lf} and \fld{freqs} columns,
the \fld{nimag} count, and the \fld{lf} cutoff property keep the raw values.
For example, a saddle point with frequencies $[1000, -300]$ still prints
\fld{nimag} 1 and \fld{lf} $-300.0$, but its Gibbs energy becomes defined
under \flg{--makereal} or \flg{--ignoreimag}.

\subsubsection*{The quasi-harmonic reference frequency}

With $\tilde{\nu}_0 = \texttt{fc} > 0$, each mode receives the free-rotor
weight

\begin{equation}
w_i = \frac{1}{1 + (\tilde{\nu}'_i/\tilde{\nu}_0)^4},
\end{equation}

and the vibrational entropy is the weighted blend

\begin{equation}
S_{\mathrm{vib}} = \sum_i \left[(1-w_i)\,S_{\mathrm{ho},i}
                                  + w_i\,S_{\mathrm{fr},i}\right].
\end{equation}

With \flg{--fc 0}, $w_i = 0$ and the entropy is purely harmonic. The
interpolation is finite for arbitrarily small positive frequencies and never
exceeds the harmonic vibrational entropy, so the quasi-harmonic Gibbs energy
is at or above the harmonic one at the same conditions.

\subsection{Species classes and rotational constants}\label{sec:species}

The geometry's mass-weighted principal moments of inertia $I_a \le I_b \le
I_c$ are computed from the coordinates and standard atomic weights.

\begin{description}
  \item[\textbf{Monatomic}] One atom. No vibrational modes and no rotational
        terms.
  \item[\textbf{Linear}] $I_a \le 10^{-4} I_c$. Two-atom geometries are
        always linear under this rule.
  \item[\textbf{Nonlinear}] Everything else with $I_b > 0$.
\end{description}

A geometry whose middle principal moment is zero or negative is degenerate:
every atom coincides with the center of mass. It has no species class and no
thermochemistry.

Rotational constants are

\begin{equation}
\tilde{B} = \frac{h}{8\pi^2 c I}
\end{equation}

in cm$^{-1}$. A linear species has one constant, $\tilde{B}_b$. A nonlinear
species has $\tilde{B}_a \ge \tilde{B}_b \ge \tilde{B}_c$, listed in
descending order by \flg{--rots}. The average moment used by the free-rotor
term is the mean of the nonzero principal moments:
$(I_a+I_b+I_c)/3$ for nonlinear species and $I$ for linear ones.

\subsection{Thermodynamic terms}\label{sec:thermo-terms}

Let $T$ be the temperature, $M$ the total mass in u, $k_B$ the Boltzmann
constant, $h$ Planck's constant, $c$ the speed of light, and

\begin{equation}
x_i = \frac{h c \tilde{\nu}'_i}{k_B T}.
\end{equation}

The terms are, in hartree per molecule:

\begin{align}
E_{\mathrm{ZPE}} &= \sum_i \tfrac{1}{2} h c \tilde{\nu}'_i / E_h, \\
E_{\mathrm{vib}} &= \sum_i \frac{h c \tilde{\nu}'_i}{e^{x_i}-1} / E_h, \\
E_{\mathrm{rot}} &=
  \begin{cases}
    0, & \text{monatomic},\\
    k_B T / E_h, & \text{linear},\\
    \tfrac{3}{2} k_B T / E_h, & \text{nonlinear},
  \end{cases}\\
E_{\mathrm{trans}} &= \tfrac{3}{2} k_B T / E_h, \\
H &= E_{\mathrm{el}} + E_{\mathrm{ZPE}} + E_{\mathrm{vib}}
     + E_{\mathrm{rot}} + E_{\mathrm{trans}} + k_B T / E_h .
\end{align}

The last term is the ideal-gas $pV = k_B T$ enthalpy contribution.

The entropies are evaluated in J~mol$^{-1}$~K$^{-1}$ with
$R = N_A k_B$:

\begin{align}
S_{\mathrm{trans}} &= R\left[\frac{5}{2}
  + \ln\!\left(\left(\frac{2\pi M u k_B T}{h^2}\right)^{3/2}
  \frac{k_B T}{p^\circ}\right)\right], \\
S_{\mathrm{rot}} &=
  \begin{cases}
    0, & \text{monatomic},\\
    R\left[\ln\!\left(\dfrac{k_B T}{\sigma h c \tilde{B}_b}\right)+1\right],
      & \text{linear},\\
    R\left[\ln\!\left(\dfrac{\sqrt{\pi}}{\sigma}
      \dfrac{(k_B T)^{3/2}}
           {\sqrt{(h c \tilde{B}_a)(h c \tilde{B}_b)(h c \tilde{B}_c)}}\right)
      +\dfrac{3}{2}\right], & \text{nonlinear},
  \end{cases} \\
S_{\mathrm{el}} &= R \ln(\text{multiplicity}).
\end{align}

Here $p^\circ = 1$ atm and $u$ is the atomic mass constant. The
quasi-harmonic vibrational entropy is

\begin{align}
S_{\mathrm{ho},i} &= R\left[
  \frac{x_i}{e^{x_i}-1} - \ln(1-e^{-x_i})\right],\\
S_{\mathrm{fr},i} &= R\left[
  \frac{1}{2} + \ln\!\left(\frac{\sqrt{8\pi^3 \mu'_i k_B T}}{h}\right)\right],\\
\mu_i &= \frac{h}{8\pi^2 c \tilde{\nu}'_i}, \qquad
\mu'_i = \frac{\mu_i \bar{I}}{\mu_i + \bar{I}},
\end{align}

where $\bar{I}$ is the average moment of inertia. The reduced moment
$\mu'_i$ caps a mode's rotor moment at the molecule's average moment. The
total entropy term and the Gibbs free energy are

\begin{equation}
T S = T\left(S_{\mathrm{trans}} + S_{\mathrm{rot}}
            + S_{\mathrm{vib}} + S_{\mathrm{el}}\right)/(N_A E_h),
\qquad
G = H - T S .
\end{equation}

The vibrational formulas are evaluated in the numerically stable
$e^{-x}$ form, so stiff modes and low temperatures underflow to zero instead
of overflowing.

\subsection{Reference values and validation}\label{sec:thermo-reference}

The derivation is validated against ORCA~6's printed quasi-RRHO values for
the fixture \texttt{1pinic1sa-110\_1\_9-str9.out} (34 atoms, 96 participating
modes, 298.15 K, reference frequency 100 cm$^{-1}$, $\sigma=1$, singlet).
The agreement, within $1\times 10^{-4}$ hartree, is:

\begin{center}
\begin{tabular}{@{}lrr@{}}
\toprule
\textbf{Quantity} & \textbf{ORCA prints (Eh)} & \textbf{CJKQC entry} \\
\midrule
Zero-point energy & 0.27547281 & \fld{zpe} \\
Thermal vibrational correction & 0.01541878 & \fld{thermal\_vibrational} \\
Thermal rotational correction & 0.00141627 & \fld{thermal\_rotational} \\
Thermal translational correction & 0.00141627 & \fld{thermal\_translational} \\
Thermal enthalpy correction ($pV$) & 0.00094421 & $k_B T/E_h$ \\
Total enthalpy & $-199.34739714$ & \fld{enthalpy} \\
Translational entropy term & 0.02035094 & \fld{ts\_translational} \\
Rotational entropy term & 0.01583464 & \fld{ts\_rotational} \\
Vibrational entropy term & 0.02738792 & \fld{ts\_vibrational} \\
Final entropy term & 0.06357350 & \fld{entropy\_term} \\
Final Gibbs free energy & $-199.41097065$ & \fld{gibbs} \\
\bottomrule
\end{tabular}
\end{center}

The rotational constants recomputed from the same geometry
(0.024157, 0.008371, 0.007528 cm$^{-1}$) match ORCA's printed values, which
also validates the moments-of-inertia path.

\subsection{A thermochemistry example}\label{sec:thermo-example}

Print the electronic energy, zero-point energy, enthalpy, entropy term, and
Gibbs free energy at 350 K with a 50 cm$^{-1}$ reference frequency, in
kcal/mol:

\begin{verbatim}
CJKQC db.jsonl --temp 350 --fc 50 --unit \
  --filename --el --zpe --h --s --g
\end{verbatim}

To make a saddle point's thermochemistry defined by treating the imaginary
mode as a real mode, add \flg{--makereal}. To drop the imaginary mode
instead, add \flg{--ignoreimag}. Only one of the two may be used. The stored
frequencies are not changed either way.

For ensemble populations derived from these Gibbs energies, see
Section~\ref{sec:populations}.

\section{Single-point corrections}\label{sec:corrections}

A structure's thermochemistry is only as good as its electronic energy. A
common workflow optimizes geometries and computes frequencies at an
affordable level, then re-evaluates the energy at a better level with separate
single-point jobs. \CJKQC{} attaches such corrections to records by name. The
stored optimization energy is never altered: the correction lives in its own
\fld{correction} entry, and corrected values appear only where you
explicitly ask for them.

There are two ways to produce the same kind of stored correction:

\begin{itemize}
  \item \flg{--add-correction} copies the energy and source details of a
        matching single-point record from another database;
  \item \flg{--extrapolate} combines same-named single points from source
        databases with a basis-set or settings extrapolation.
\end{itemize}

Everything that follows is shared: the \fld{elcorr}, \fld{hcorr}, and
\fld{gcorr} columns, the \fld{elcorr} and \fld{gcorr} cutoff properties and
sort properties, the corrected formation and solvation sections, and the
corrected population levels.

\subsection{The correction database}\label{sec:correction-database}

A correction database is an ordinary \CJKQC{} database, usually built by a
normal collection run over the single-point outputs. Each correction record
therefore carries its own source information: \fld{program},
\fld{program_version},
and \fld{command}. A database assembled by merging or by hand works the same
way; whatever a record lacks is simply not copied.

A \emph{candidate} is a correction record with status \fld{ok}, a present
\fld{electronic_energy}, and a present \fld{name}. Only candidates attach.
Two candidates sharing one name make the database ambiguous: that is a fatal
error naming the name, and nothing is written.

\subsection{Normal corrections with \texttt{-{}-add-correction}}\label{sec:add-correction}

\begin{verbatim}
CJKQC optimized.jsonl --add-correction single-points.jsonl \
  --filename --el --elcorr --g --gcorr -o corrected.jsonl
\end{verbatim}

The flag takes one path, is a manipulation flag, and cannot be repeated. The
path is always loaded as a database, never classified like the positional
inputs. A missing or unreadable file is fatal. A corrupt line is skipped and
reported, and the valid records still attach.

Matching is exact, case-sensitive equality of \fld{name}. Every input record
whose name equals a candidate's name receives the correction, including
several records sharing that name. Input records without a name never match.
Attachment is independent of the input record's status; the usual rule that
only \fld{ok} records are shown decides visibility.

The stored correction copies the candidate's values:

\begin{itemize}
  \item \fld{energy}: the candidate's \fld{electronic_energy}, in hartree;
  \item \fld{program} and \fld{program_version}: copied from the candidate;
  \item \fld{commands}: a one-element list holding the candidate's
        \fld{command}, or absent when it has none;
  \item \fld{kind}: \texttt{sp} for an attached single-point record.
\end{itemize}

A record that already carries a correction has it replaced by a new match. A
record the correction database does not name keeps whatever correction it
already stores, so re-running a corrected database against a larger
correction database only adds and updates. The changes apply only to the
current run; \flg{--out} writes the corrected records. Without \flg{--out},
nothing is stored.

Warnings report candidates skipped for status or missing energy, unmatched
candidates, and corrupt lines. Warnings do not change the exit code or the
output, except that a correction database that loads no record at all while
reporting load problems is fatal.

\subsection{Corrected values}\label{sec:corrected-values}

Three corrected values become available:

\begin{description}
  \item[\fld{elcorr}] The attached correction's energy, in hartree.
        Undefined without a correction.
  \item[\fld{hcorr}] The enthalpy derived at the run's conditions from the
        correction energy and the record's other stored values: the standard
        derivation with the correction energy in place of
        \fld{electronic_energy}.
  \item[\fld{gcorr}] The Gibbs free energy derived the same way.
\end{description}

The raw levels are untouched. \fld{el}, \fld{g}, \fld{zpe}, \fld{h},
\fld{s}, the cutoffs and sorts on them, \flg{--glob}, \flg{--uniq}, and the
raw formation section all keep using the optimization-level values.
Corrected values appear only where requested.

\begin{note}
A record loaded from a previously corrected database carries its correction
with it. No flag is needed to use stored corrections; \flg{--add-correction}
is only needed to add or update them.
\end{note}

The corrected print columns are \flg{--elcorr}, \flg{--hcorr}, and
\flg{--gcorr}. \flg{--elcorr} prints the correction energy itself;
\flg{--hcorr} and \flg{--gcorr} print the corrected enthalpy and Gibbs
energy. \flg{--correction} prints the source details, for example:

\begin{verbatim}
sp orca 6.0.0 "! wB97x-3c def2-TZVPP Sp"
\end{verbatim}

\subsection{Extrapolation with \texttt{-{}-extrapolate}}\label{sec:extrapolation}

\flg{--extrapolate} combines same-named single points from source databases
instead of copying one energy. The mode word leads; the remaining values are
the source paths and the mode's parameters, in JKextrapolate's order.

\begin{center}
\begin{tabular}{@{}L{4.4cm}L{9.1cm}@{}}
\toprule
\textbf{Invocation} & \textbf{Meaning} \\
\midrule
\flg{--extrapolate 2p LOW HIGH N ALPHA BETA} &
Two-point basis-set extrapolation to the complete-basis-set limit. LOW holds
the lower cardinal number's single points, HIGH the next one. $N$ is the
lower cardinal number (a whole number $\ge 1$), and $\alpha>0$, $\beta>0$ are the
SCF and correlation exponents. Produces kind \texttt{cbs}. \\
\flg{--extrapolate 2s LOW HIGH F} &
A CPS/LAF settings blend over two databases that differ in calculation
settings at fixed basis. $F$ is the blend fraction. Produces kind
\texttt{cps}. \\
\flg{--extrapolate 2p LOWPAIR HIGHPAIR N ALPHA BETA F} &
The combined form. LOWPAIR and HIGHPAIR are comma-joined path pairs. The
first entry of each pair is the low-settings calculation and the second is
the high-settings calculation. One two-point extrapolation runs per settings
level, then the two results blend. Produces kind \texttt{cbs-cps}. \\
\bottomrule
\end{tabular}
\end{center}

The two \texttt{2p} forms share the mode word because both perform a two-point
basis-set extrapolation. The combined form is recognized by its arguments: it
takes six values instead of five, and its two source arguments are
comma-joined path pairs.

Every value must be valid. An unknown mode, a wrong value count, a non-numeric
or non-whole-number parameter, $N<1$, $\alpha$ or $\beta$ at or below zero, a
non-finite $F$, a comma group with more than two entries, a comma group where
a single path is required, or a trailing $F$ without two comma-joined pairs
is a usage error. Usage errors happen before any input is read.

\begin{warning}
Like \flg{--pop}, \flg{--extrapolate} eats the values that follow it until
the next \flg{--flag}. Place it after the positional inputs, or a positional
path after it may be swallowed as a parameter. Only one \flg{--extrapolate}
is allowed per run, and it cannot be combined with \flg{--add-correction}
because a record carries one correction.
\end{warning}

\subsubsection*{The formulas}

All energies are in hartree. For a two-point extrapolation over points $X$
(low basis, cardinal $N$) and $Y$ (high basis, cardinal $N+1$), let

\begin{equation}
e_X = e^{-\alpha\sqrt{N}}, \qquad
e_Y = e^{-\alpha\sqrt{N+1}}.
\end{equation}

Then

\begin{equation}
E_{\mathrm{CBS}} =
\frac{E_{\mathrm{scf},X}\,e_Y - E_{\mathrm{scf},Y}\,e_X}{e_Y-e_X}
+
\frac{N^\beta E_{\mathrm{corr},X} - (N+1)^\beta E_{\mathrm{corr},Y}}
     {N^\beta - (N+1)^\beta},
\end{equation}

where the SCF term extrapolates the records' \fld{scf_energy} and the
correlation term extrapolates the residual
\fld{electronic_energy} $-$ \fld{scf_energy}. For a DFT record the residual
is zero and the SCF term carries the whole energy. A record whose stored
electronic energy is itself the SCF value has a zero residual too, so its
extrapolation is an SCF-limit energy, not a correlated one. Extrapolation
never recovers a correlation energy that the record does not store.

The settings blend is

\begin{equation}
E = E_{\mathrm{low}} + F\,(E_{\mathrm{high}} - E_{\mathrm{low}}),
\end{equation}

over the records' electronic energies. In the combined form, let
$E_{\mathrm{CBS,low}}$ be the two-point result for the first entries of both
pairs and $E_{\mathrm{CBS,high}}$ the result for the second entries. Then

\begin{equation}
E = E_{\mathrm{CBS,low}}
  + F\,(E_{\mathrm{CBS,high}} - E_{\mathrm{CBS,low}}).
\end{equation}

A non-finite result leaves the record without a correction and is warned like
any other skip; an infinite or invalid energy is never written.

\subsubsection*{Recommended basis-set coefficients}

The SCF form and the coefficients below follow Neese and co-workers (2011);
the correlation form follows Helgaker and co-workers (1997). The two-point
form needs a pair of coefficients for the basis family and the cardinal
numbers being combined. The table lists recommended values. The numbers in
parentheses are the cardinal numbers; use the pair that matches the lower
basis.

\begin{center}
\begin{tabular}{@{}lrrrr@{}}
\toprule
Basis set & $\alpha(2,3)$ & $\beta(2,3)$ & $\alpha(3,4)$ & $\beta(3,4)$ \\
\midrule
\texttt{cc-pVXZ}      & 4.42 & 2.46 & 5.46 & 3.05 \\
\texttt{aug'-cc-pVXZ} & 4.30 & 2.51 & 5.79 & 3.05 \\
\texttt{aug-cc-pVXZ}  & 4.30 & 2.51 & 5.79 & 3.05 \\
\bottomrule
\end{tabular}
\end{center}

The reduced \texttt{aug'} sets use the fully augmented coefficients: no
separate fit exists for them, and the two families are expected to give
similar coefficients.

For example, extrapolating triple- and quadruple-zeta energies of an
augmented basis set uses $\alpha=5.79$ and $\beta=3.05$:

\begin{verbatim}
CJKQC run.jsonl --extrapolate 2p cbs/tz.jsonl cbs/qz.jsonl 3 5.79 3.05 \
  -o cbs.jsonl
\end{verbatim}

\subsubsection*{Settings extrapolations: CPS, LAF, and PNO}

Basis-set extrapolation raises the cardinal number; CPS and LAF instead
combine two calculations that differ in accuracy settings at a fixed basis.
Both use the same blend with $F=1.5$:

\begin{equation}
E = E_{\mathrm{low}} + 1.5\,(E_{\mathrm{high}} - E_{\mathrm{low}}).
\end{equation}

To run one, build a database for each settings level and pass the two
databases to the \texttt{2s} form, low level first:

\begin{verbatim}
CJKQC run.jsonl --extrapolate 2s cps/pno3.jsonl cps/pno4.jsonl 1.5 \
  -o cps.jsonl
\end{verbatim}

For a DLPNO calculation, a PNO threshold of $3.33\times10^{-T}$ is written
PNO($T$), so the command above blends PNO(3) and PNO(4) and is called
CPS(3,4). For an LNO calculation, use two successive locality settings, for
example vLoose and Loose; the shortened names are vL and L, so that pair is
called LAF(vL,L). When basis and settings change together, use the combined
\texttt{2p} form: give two comma-joined path pairs (low and high settings at
the lower basis, then the same two settings at the higher basis) and add
$F=1.5$ after the coefficients. \CJKQC{} extrapolates each pair to the
basis-set limit and then blends the two results with the same formula.

For example, to combine triple- and quadruple-zeta calculations at the
vLoose and Loose LNO settings, give the lower basis at vLoose and Loose as
the first pair and the higher basis in the same settings order as the second
pair:

\begin{verbatim}
CJKQC run.jsonl --extrapolate 2p \
  cbs/tz-vl.jsonl,cbs/tz-l.jsonl cbs/qz-vl.jsonl,cbs/qz-l.jsonl \
  3 5.46 3.05 1.5 -o lno-cbs.jsonl
\end{verbatim}

This performs a CBS(3,4) extrapolation at each LNO setting and then blends
the two limits, producing kind \texttt{cbs-cps}.

\subsubsection*{Candidate and matching rules}

A source candidate is a source-database record with status \fld{ok}, a
present \fld{name}, a present \fld{electronic_energy}, and, for every mode
that extrapolates over bases, a present \fld{scf_energy}. A source record
failing any of these is skipped and warned individually.

The following conditions are fatal and write nothing:

\begin{itemize}
  \item two candidates with equal names within one source database;
  \item contributors of one name with different present \fld{program}
        values;
  \item contributors of one name with different present
        \fld{program_version} values;
  \item contributors of one name with exactly equal \fld{electronic_energy}
        values, which means the same calculation entered the sources twice.
\end{itemize}

The program and version checks apply whether or not a run record names the
name; only present-and-different values violate. The duplicate-energy check
covers electronic energies only: within-basis pairs in the combined form may
legitimately share the SCF energy.

A run record receives the extrapolated correction if and only if its name
appears as a candidate in \emph{every} source database of the invocation.
Matching is exact and case-sensitive. Every run record sharing a fully
covered name receives the correction. A run record whose name is not fully
covered is skipped with a warning naming the name and the first source
database missing it. A record that already carries a correction has it
replaced on a match; an uncovered record keeps its stored correction.

\subsubsection*{The attached extrapolation correction}

The stored correction contains:

\begin{itemize}
  \item \fld{energy}: the mode's result in hartree;
  \item \fld{kind}: \texttt{cbs}, \texttt{cps}, or \texttt{cbs-cps};
  \item \fld{commands}: the contributing records' commands in contributor
        order, omitting records without a command;
  \item \fld{program} and \fld{program_version}: the last contributor's
        values.
\end{itemize}


\section{Formation energies}\label{sec:formation}

Clusters are put in context by their formation energies: how much a
structure's energy differs from a stoichiometric mix of isolated,
single-component reference structures. A cluster database already contains
its own references, namely the monomers. \CJKQC{} scans the run's final
records for monomers and prints each printed record's formation energy and
formation free energy as a section after the main table.

Nothing is stored. Formation values are derived at request time from the
records' stored values under the run's thermochemistry conditions, so
temperature, reference frequency, and unit changes apply to them like to any
other derived property.

The sign convention is:

\begin{quote}
\textbf{negative} = more stable than the isolated monomers.
\end{quote}

The binding energy is the same number with the opposite sign; there is no
separate binding column.

\subsection{Monomer references}\label{sec:monomer-references}

A \textbf{monomer record} is an \fld{ok} record whose composition is exactly
one component with count 1. Thus \texttt{\{"am": 1\}} qualifies, while
\texttt{\{"am": 2\}} and \texttt{\{"am": 1, "sa": 1\}} do not. Records
without components never qualify.

Monomers are collected from exactly the records the run prints: the records
that remain after every filtering step in a filtering run, or every loaded
record in a print-only run. A monomer that a filter removes is gone:
it prints no row and serves as no reference.

For each component code, the reference is a single monomer record: the first
of that code's monomers under the run's sort ordering.

\begin{center}
\begin{tabular}{@{}ll@{}}
\toprule
\textbf{Sort} & \textbf{Reference monomer} \\
\midrule
\fld{el}, or no sort & lowest electronic energy \\
\fld{g} & lowest Gibbs free energy at the run's conditions \\
\fld{rg} & first under radius of gyration \\
\fld{mass} & first under total mass \\
\fld{filename} & first under the file-stem ordering (nameless first) \\
\bottomrule
\end{tabular}
\end{center}

The ordering is always ascending; \flg{--reverse} orients the table, never
the reference. Ties keep the first in load order. For numeric properties, a
monomer whose sorting value is undefined orders after every defined one.

Both reference values come from that one record:

\begin{itemize}
  \item the \textbf{electronic reference} is its \fld{electronic\_energy};
  \item the \textbf{Gibbs reference} is its Gibbs free energy derived at the
        run's conditions.
\end{itemize}

Under \flg{--sort g}, the lowest-G conformer supplies the electronic
reference too, even when another conformer's electronic energy is lower. A
level whose value the chosen record lacks has no reference at that level. A
code with no qualifying monomer has no reference.

\subsection{The raw formation section}\label{sec:raw-formation}

\flg{--formation} is a print flag: it produces output and satisfies the
visible-outcome rule on its own, but adds no column to the normal table. The
section prints after the main table (if one was requested) and before any
summary line.

For a printed record with components $C = \{\text{code}_i : n_i\}$:

\begin{align}
E_{\mathrm{form}} &= E_{\mathrm{el}}
  - \sum_i n_i\,E_{\mathrm{ref}}(\text{code}_i),\\
G_{\mathrm{form}} &= G - \sum_i n_i\,G_{\mathrm{ref}}(\text{code}_i).
\end{align}

Both reference values are the ones selected above. A level is undefined when
the record lacks \fld{components}, lacks the level's input energy (electronic
energy for the electronic level, defined thermochemistry for the Gibbs
level), or any referenced code lacks that level's monomer reference.
Undefined is a property of the level: a record can have a defined formation
energy and an undefined formation free energy, and vice versa.

The section is a table with header \fld{name el g}, one row per printed
record in run order. \fld{name} is the short composition label (for example
\texttt{2am3sa}), or \texttt{-} when the record has no components.
\fld{el} and \fld{g} are the formation energy and formation free energy in
the run's energy unit: hartree by default, kcal/mol under \flg{--unit}. An
undefined value prints \texttt{-}.

\subsection{The corrected formation section}\label{sec:corrected-formation}

\flg{--formation-corr} prints a second section evaluated at the attached
single-point corrections. Every input energy is a corrected value; a record
with no correction has no values at all, so corrected and raw levels are
never mixed.

The reference monomers are the same set as for the raw section, selected the
same way. A code's corrected electronic reference is that monomer's
correction energy, and its corrected Gibbs reference is the monomer's
\fld{gcorr} value. A level's reference is undefined when the selected
monomer carries no correction, or its corrected Gibbs is undefined.

For a printed record:

\begin{align}
E_{\mathrm{form,corr}} &= E_{\mathrm{corr}}
  - \sum_i n_i\,E_{\mathrm{corr,ref}}(\text{code}_i),\\
G_{\mathrm{form,corr}} &= G_{\mathrm{corr}}
  - \sum_i n_i\,G_{\mathrm{corr,ref}}(\text{code}_i).
\end{align}

The section has header \fld{name elcorr gcorr}. It prints after the raw
formation section, when one was requested. A run with
\flg{--formation-corr} but no other print flag prints only the corrected
section.

\subsection{Concentrations and \texttt{-{}-cnt}}\label{sec:concentrations}

\flg{--conc CODE VALUE} evaluates the formation free energies at a vapor
concentration different from the standard state. It may be repeated, with
one concentration setting per code. The code must be a non-empty, letter-only
component code. When a formation section is printed, every such code must
appear as a component of at least one record in the section's record set.
Specifying the same code twice is a usage error.

VALUE is a partial pressure in pascals, optionally suffixed:

\begin{center}
\begin{tabular}{@{}lll@{}}
\toprule
\textbf{Form} & \textbf{Meaning} & \textbf{Pressure in Pa} \\
\midrule
\texttt{100} & pascals & $100$ \\
\texttt{100ppt} & parts per trillion by mole fraction & $100\cdot10^{-12}p^\circ$ \\
\texttt{10ppb} & parts per billion by mole fraction & $10\cdot10^{-9}p^\circ$ \\
\texttt{1e10cmmc} & molecules per cm$^3$ & $c\cdot10^{6}k_B T$ \\
\bottomrule
\end{tabular}
\end{center}

Here $p^\circ = 101325$ Pa. The value must be finite and positive; the
\texttt{\^{}} power syntax is not part of the grammar. For powers of ten, use
scientific notation instead, for example \texttt{1e10} or \texttt{1E10} for
$10^{10}$. The concentration term
is a chemical-potential (entropy) correction, so it shifts only the Gibbs
levels: the raw section's \fld{g} and the corrected section's \fld{gcorr}.
The electronic levels \fld{el} and \fld{elcorr} stay at the standard state.

For the set of concentration settings $S$ and a record with components $C$:

\begin{equation}
g_{\mathrm{conc}} = g_{\mathrm{form}}
  - \sum_{\text{code}_i \in S} n_{i,\mathrm{eff}}\,
    \frac{R T \ln(c_i/p^\circ)}{N_A E_h},
\end{equation}

where $c_i$ is the resolved partial pressure and $n_{i,\mathrm{eff}} = n_i$.
Under \flg{--cnt}, the self-consistent form is used:

\begin{equation}
n_{i,\mathrm{eff}} = n_i\,\frac{N-1}{N},
\qquad N = \sum_i n_i .
\end{equation}

The self-consistent form keeps the concentration term of a monomer at zero
under any concentration, so a monomer's formation free energy stays exactly
zero while a non-reference conformer keeps its conformer difference. A
concentration setting whose code is absent from the record's components
contributes nothing.

\subsection{Boltzmann-averaged sections}\label{sec:bavg}

\flg{--bavg} collapses each printed formation section into one
Boltzmann-averaged row per composition. It is a settings flag: it configures
the section's presentation, never counts as a manipulation or print flag, and
is an accepted no-op without a printed formation section. The main table is
never affected; it still prints one row per record.

The collapse groups the section's record set by \fld{components}, always by
composition (\flg{--noname} never pools the section), in first-appearance
order. Records without components form one unnamed ensemble. A group's
members are its \fld{ok} records that define the section's Gibbs level: the
derived Gibbs for the raw section, \fld{gcorr} for the corrected one. Failed
and truncated records contribute nothing, and an \fld{ok} record without the
level value contributes nothing.

An ensemble with members carries two values:

\begin{itemize}
  \item the \textbf{electronic value} is the members' population-weighted
        mean $\sum_j p_j e_j$, with the members' Boltzmann populations at the
        run's temperature. Here $e_j$ is \fld{electronic\_energy} for the raw
        section and the correction energy for the corrected section;
  \item the \textbf{Gibbs value} is the ensemble free energy
        \begin{equation}
        G_{\mathrm{ens}} = -k_B T \ln\!\left(\sum_j e^{-g_j/(k_B T)}\right)
        = g_{\min} - k_B T \ln Z,
        \end{equation}
        computed against the members' minimum as in the population
        numerics. This keeps the conformational entropy term
        $-k_B T\ln Z$; a plain weighted mean would drop it.
\end{itemize}

The two values are defined together or not at all: a member that defines the
Gibbs level always defines the electronic level, because the Gibbs derives
from the electronic energy.

Over the collapsed set, the standing rules apply unchanged. An ensemble whose
composition is exactly one component with count 1 is its code's monomer
ensemble and supplies both reference values; a monomer row's formation values
are therefore exactly zero. An ensemble without values supplies no reference.

The section prints one row per ensemble in first-appearance order, using the
same header, unit, and formatting as the uncollapsed section. A composition
whose records all lack the section's Gibbs level prints one row of dashes.
Undefined members never poison their group: the surviving members define the
ensemble.

\subsection{A formation example}\label{sec:formation-example}

Print formation energies and free energies for the lowest conformer of each
composition, using the lowest-Gibbs monomer of each code as the reference:

\begin{verbatim}
CJKQC db.jsonl --sort g --glob --formation --unit \
  --filename --name --g
\end{verbatim}

Evaluate the same formation free energies at an ammonia vapor of
$10^{-9}$ bar-equivalent mole fraction, with the self-consistent form:

\begin{verbatim}
CJKQC db.jsonl --formation --conc am 10ppb --cnt --unit --name
\end{verbatim}

Collapse the section into one Boltzmann-averaged row per composition:

\begin{verbatim}
CJKQC db.jsonl --formation --bavg --unit --name
\end{verbatim}

In the last example the main table, if requested, would still show one row
per record; only the formation section is collapsed.

\section{Solvation and hydration distributions}\label{sec:solvation}

Clusters in equilibrium with a vapor exchange solvent molecules. At a given
relative humidity, the hydrates of one core composition populate a Boltzmann
distribution over the solvent count. \CJKQC{} prints this distribution with
\flg{--hydration} for water and \flg{--solvation CODE} for any other solvent.

The distribution is built on the formation machinery of
Chapter~\ref{sec:formation}. A state's distribution value is its formation
free energy with the solvent's chemical potential set to its vapor pressure.
The distribution at a given relative humidity and the formation energy at a
given concentration are therefore two views of the same calculation.

\subsection{Flags and solvent pressure}\label{sec:solvation-flags}

\begin{center}
\begin{tabular}{@{}L{4.1cm}L{9.4cm}@{}}
\toprule
\textbf{Flag} & \textbf{Meaning} \\
\midrule
\flg{--hydration} & Print the raw distribution section for water, solvent
code \texttt{w}. \\
\flg{--hydration-corr} & Print the corrected distribution section for water;
requires corrections, attached or stored. \\
\flg{--solvation CODE} & Print the raw distribution section for the named
solvent code. \\
\flg{--solvation-corr CODE} & Print the corrected distribution section for
the named solvent code. \\
\flg{--rh VALUE} & Relative humidity as a fraction: \texttt{0.5} means 50
\%. Water only. \\
\flg{--psolvent VALUE} & Solvent partial pressure, using the \flg{--conc}
VALUE grammar. Valid for every solvent and overrides \flg{--rh}. \\
\bottomrule
\end{tabular}
\end{center}

All four section flags are print flags and each is a visible outcome on its
own. The plain spellings print the raw-level section; the \texttt{-corr}
spellings print the corrected-level section. The raw section prints before the
corrected one, whatever order the flags were given.

All solvent flags in one run must agree on a single solvent code. A second,
different code is a usage error, while repeated spellings of the same code are
redundant. Codes are non-empty, letter-only, case-significant strings.

\flg{--rh} is water physics: it must be a finite fraction greater than zero,
and values above 1 (supersaturation) are accepted. It is a usage error with a
solvent code other than \texttt{w}. \flg{--psolvent} accepts the same values
as \flg{--conc}: pascals, or a value suffixed with \texttt{ppt}, \texttt{ppb},
or \texttt{cmmc}.

The pressure resolution order is:

\begin{enumerate}
  \item \flg{--psolvent}, if given;
  \item otherwise \flg{--rh}, if given;
  \item otherwise the saturated vapor pressure, which is valid for water only.
\end{enumerate}

For a non-water solvent, neither \flg{--rh} nor the default applies: an
explicit \flg{--psolvent} is required.

For water, \flg{--rh} uses the Magnus correlation

\begin{equation}
p_{\mathrm{sat}}(T) = 610.94\,
  \exp\!\left(\frac{17.625\,\tau}{243.04+\tau}\right)\ \mathrm{Pa},
\qquad \tau = T - 273.15,
\end{equation}

which gives 610.94 Pa at 273.15 K and about $3.162\times10^3$ Pa at 298.15 K.
The resolved pressure $p_w = \texttt{rh}\cdot p_{\mathrm{sat}}(T)$ must be a
finite positive number before the run starts; a temperature at the domain edge
that makes it non-finite is a usage error.

The resolved pressure enters the run as a concentration setting supplied
automatically, as if \flg{--conc CODE p} had been given. The formation
sections' solvent term shifts accordingly, \flg{--cnt} scales it like any
concentration setting, and the distribution values carry it. Naming the
solvent code in an explicit \flg{--conc} is a usage error, because the
automatic setting already owns it; settings for other components ride along
untouched.

\subsection{States, cores, and distribution values}\label{sec:solvation-states}

The sections evaluate exactly the record set the formation sections evaluate:
the records that remain after filtering, every loaded record for a print-only
run, or the collapsed ensembles under \flg{--bavg}.

A \textbf{state} is one composition. Its \textbf{core} is that composition
without the solvent code. For example, with solvent \texttt{w}, the state
\texttt{1sa2w} has core \texttt{1sa}, and the states \texttt{1sa},
\texttt{1sa1w}, and \texttt{1sa2w} belong to the same core group. Records and
ensembles without components are excluded. States whose core is empty, such
as the pure solvent clusters \texttt{1w}, \texttt{2w}, and so on, are also
excluded: the solvent monomer itself prints no distribution row.

Core groups print in first-appearance order of their core among the printed
set. States within a group print in first-appearance order of their
composition.

A state's distribution value at the section's level is:

\begin{itemize}
  \item \textbf{uncollapsed}: the lowest free conformer, that is, the minimum
        over the state's members' defined level values. Members with an
        undefined value are ignored; a state with no defined member is
        undefined;
  \item \textbf{under \flg{--bavg}}: the ensemble's level value, because the
        collapse already provides one ensemble per composition. The
        collapse replaces the conformer minimum, including conformational
        entropy.
\end{itemize}

A state's formation value needs the monomer references of its full
stoichiometry. A missing or level-less reference code leaves the state
undefined at that level. A hydrated state whose solvent monomer was filtered
away prints a dash, while the solvent-free states of the same core can remain
defined.

\subsection{The distribution}\label{sec:distribution}

Within a core group, each defined state's weight is

\begin{equation}
w_{\mathrm{state}} = \exp\!\left(
  -\frac{v_{\mathrm{state}} - v_{\min}}{k_B T}\right),
\end{equation}

where $v_{\min}$ is the group's minimum defined value and $T$ is the run's
temperature. Its percentage is

\begin{equation}
100\,\frac{w_{\mathrm{state}}}{\sum_{\mathrm{defined}} w}.
\end{equation}

Undefined states contribute nothing and print \texttt{-}. A group with no
defined states prints all dashes. A group with exactly one defined state
prints 100.0.

The pressure physics rides inside the values. The implied solvent setting's term
$-n R T \ln(p_w/p^\circ)$ makes the weights equal the JKQC form
$(p_w/p^\circ)^n \exp(-(G-nG^\circ_{\mathrm{solvent}})/(k_B T))$. Every
non-solvent term common to the group cancels in the normalization.
\flg{--cnt}'s $(N-1)/N$ scaling varies with the state's total molecule count,
so under \flg{--cnt} the distribution simply follows the formation values.
The minimum subtraction keeps every exponent at or below zero, so no overflow
is reachable and a state far above the minimum underflows to zero without
affecting the rest.

\flg{--noname} never pools the core groups; the grouping is intrinsic to the
distribution. Sorting and \flg{--reverse} reorder the main table, never the
section, which keeps its own grouping order.

\subsection{Printed layout}\label{sec:solvation-layout}

There is no header row. For each core group, two rows are printed: the
states' composition tokens, then their percentages. Percentages print with
one decimal, \texttt{-} for an undefined state, and need not sum to exactly
100.0 because of rounding. A blank line follows every group's percentage row,
so groups read as blocks.

Columns align within a group: each column is as wide as the wider of its label
and percentage cell, cells are right-aligned, and there is a single space
between columns. For example:

\begin{verbatim}
 1sa 1sa1w 1sa2w
40.0  45.0  15.0

 1pinic 1pinic1w
 72.3    27.7
\end{verbatim}

The standard-output order is: main table, raw formation section, corrected
formation section, hydration/solvation sections (raw before corrected), then
the completion summary. The sections print over the same record set as the
formation sections, whether or not any formation section was requested.

\subsection{Solvation examples}\label{sec:solvation-example}

Print the water distribution at 50 \% relative humidity:

\begin{verbatim}
CJKQC hydrated.jsonl --hydration --rh 0.5 --filename --name
\end{verbatim}

Print both raw and corrected water distributions at 298.15 K:

\begin{verbatim}
CJKQC corrected.jsonl --hydration --hydration-corr --temp 298.15
\end{verbatim}

Print a sulfuric-acid solvation distribution at an explicit partial pressure
of $10^{10}$ molecules per cm$^3$:

\begin{verbatim}
CJKQC clusters.jsonl --solvation sa --psolvent 1e10cmmc
\end{verbatim}

Combine a concentration for another component with the water distribution:

\begin{verbatim}
CJKQC clusters.jsonl --hydration --rh 0.8 --conc am 10ppb
\end{verbatim}

Here the \flg{--conc am ...} setting shifts the formation values of the
ammonia component, while the automatic water setting carries the humidity; the
distribution normalizes within each core group.


\appendix
\section{Complete flag reference}\label{app:flags}

This appendix lists every flag accepted by \CJKQC{} 0.1.0. Flags are grouped
by purpose; within a group the order is not significant. Words in angle
brackets are placeholders you replace with your own value, and square brackets
mark optional values.

\subsection{Inputs and output}

\begin{longtable}{@{}L{4.0cm}L{9.5cm}@{}}
\toprule
\textbf{Flag} & \textbf{Meaning} \\
\midrule
\endhead
\texttt{INPUTS...} & Positional, required. Program outputs (build mode) or
record databases (merge mode). \\
\flg{-o}, \flg{--out <PATH>} & Database to write. Required unless a
manipulation or print flag makes it optional. \\
\flg{--fout} & Force replacement of an existing \flg{--out} or
\texttt{movie.xyz} without asking. \\
\flg{--head-bytes <N>} & How much of the start of each program output is read,
in bytes (default 262144, about 256 kilobytes). \\
\flg{--tail-bytes <N>} & How much of the end of each program output is read,
in bytes (default 4194304, about 4 megabytes). \\
\flg{--parallel <N>} & How many files are read at once (default 1). Merge
mode ignores it. \\
\flg{--movie} & Write final structures to \texttt{movie.xyz}, one frame per
\fld{ok} record with a geometry. \\
\bottomrule
\end{longtable}

\subsection{Print columns}

There is no default column selection. Columns always appear in the fixed
order of Section~\ref{sec:printing}.

\begin{longtable}{@{}L{3.5cm}L{10.0cm}@{}}
\toprule
\textbf{Flag} & \textbf{Column and quantity} \\
\midrule
\endhead
\flg{--filename} & \fld{filename}: the record's \fld{name} (file stem). \\
\flg{--name} & \fld{name}: composition written without separators
(\texttt{2am3sa}). \\
\flg{--path} & \fld{path}: the full path of the original output file. \\
\flg{--status} & \fld{status}: \fld{ok}, \fld{failed}, or \fld{truncated}. \\
\flg{--program} & \fld{program}: program name and version (\texttt{xtb 6.7.1}). \\
\flg{--command} & \fld{command}: the run's command line, double-quoted. \\
\flg{--correction} & \fld{correction}: correction kind and source details. \\
\flg{--el} & \fld{el}: electronic energy. \\
\flg{--zpe} & \fld{zpe}: harmonic zero-point energy. \\
\flg{--h} & \fld{h}: enthalpy at \flg{--temp}. \\
\flg{--s} & \fld{s}: entropy term $T\,S$ at \flg{--temp}. \\
\flg{--g} & \fld{g}: Gibbs free energy at \flg{--temp}. \\
\flg{--elcorr} & \fld{elcorr}: attached correction energy. \\
\flg{--hcorr} & \fld{hcorr}: enthalpy from the correction energy. \\
\flg{--gcorr} & \fld{gcorr}: Gibbs free energy from the correction energy. \\
\flg{--pop [LEVEL]} & Population column; \texttt{LEVEL} is \fld{g} (default),
\fld{el}, \fld{elcorr}, or \fld{gcorr}. Repeatable. \\
\flg{--rg} & \fld{rg}: mass-weighted radius of gyration in \AA{}. \\
\flg{--mass} & \fld{mass}: total mass in u. \\
\flg{--dip} & \fld{dip}: dipole magnitude in atomic units. \\
\flg{--dips} & \fld{dips}: dipole vector components in atomic units. \\
\flg{--pol} & \fld{pol}: isotropic polarizability in atomic units. \\
\flg{--pols} & \fld{pols}: polarizability tensor values in atomic units. \\
\flg{--charge} & \fld{charge}: net charge. \\
\flg{--multiplicity} & \fld{multiplicity}: spin multiplicity. \\
\flg{--rsn} & \fld{rsn}: rotational symmetry number. \\
\flg{--natoms} & \fld{natoms}: geometry atom count. \\
\flg{--lf} & \fld{lf}: lowest stored frequency in cm$^{-1}$. \\
\flg{--freqs} & \fld{freqs}: all stored frequencies in cm$^{-1}$. \\
\flg{--nfreq} & \fld{nfreq}: number of stored frequencies. \\
\flg{--nimag} & \fld{nimag}: number of negative frequencies. \\
\flg{--rots} & \fld{rots}: rotational constants in cm$^{-1}$. \\
\flg{--components} & \fld{components}: composition as \texttt{2am+3sa}. \\
\bottomrule
\end{longtable}

\subsection{Filtering and manipulation}

\begin{longtable}{@{}L{4.3cm}L{9.2cm}@{}}
\toprule
\textbf{Flag} & \textbf{Meaning} \\
\midrule
\endhead
\flg{--extract <QUERY>} & Keep records whose composition matches the query. \\
\flg{--nextract <QUERY>} & Drop records whose composition matches the query. \\
\flg{--rename <OLD> <NEW>} & Move the count of component code \texttt{OLD} to
code \texttt{NEW}. Repeatable; renames apply in the order given. \\
\flg{--cut <PROP> <V>} & Absolute cutoff, property $\le V$. \\
\flg{--filter-lt <PROP> <V>} & Absolute cutoff, property $< V$. \\
\flg{--pass <PROP> <V>} & Absolute cutoff, property $> V$. \\
\flg{--filter-ge <PROP> <V>} & Absolute cutoff, property $\ge V$. \\
\flg{--is <PROP> <V>} & Absolute cutoff, property $= V$. \\
\flg{--isnot <PROP> <V>} & Absolute cutoff, property $\ne V$. \\
\flg{--cutr <PROP> <V>} & Relative cutoff, property $\le V$ above the group
minimum. \\
\flg{--rel-filter-lt <PROP> <V>} & Relative cutoff, $<V$ above the minimum. \\
\flg{--passr <PROP> <V>} & Relative cutoff, $>V$ above the minimum. \\
\flg{--rel-filter-ge <PROP> <V>} & Relative cutoff, $\ge V$ above the
minimum. \\
\flg{--uniq <MODE>} & \texttt{dup}, or a comma list of
\fld{metric[decimals]} properties. \\
\flg{--reacted} & Keep the composition group's majority fragment pattern. \\
\flg{--reacted2} & Keep only the deviating fragment patterns. \\
\flg{--sort <PROPERTY>} & Sort by \fld{el}, \fld{g}, \fld{elcorr}, \fld{gcorr},
\fld{rg}, \fld{mass}, or \fld{filename}. \\
\flg{--reverse} & Reverse the sort orientation. \\
\flg{--select <N>} & Keep the first $N$ records of each group after sorting. \\
\flg{--glob} & Keep each group's lowest-Gibbs record. \\
\flg{--shuffle} & Put the final result in a random order; the same seed gives
the same order. \\
\flg{--seed <N>} & Seed for \flg{--shuffle} (default 0). \\
\flg{--noname} & Treat the whole database as one group; \flg{--extract}
compares names exactly. \\
\bottomrule
\end{longtable}

\begin{note}
The help text for \flg{--sort} mentions ``name'', but the accepted value is
\fld{filename}, which sorts by the record's \fld{name}. Use
\flg{--sort filename}.
\end{note}

\subsection{Thermodynamics}

\begin{longtable}{@{}L{3.5cm}L{10.0cm}@{}}
\toprule
\textbf{Flag} & \textbf{Meaning} \\
\midrule
\endhead
\flg{--temp <K>} & Temperature in K (default 298.15; must be $>0$). \\
\flg{--fc <CM-1>} & Quasi-harmonic reference frequency (default 100; 0 selects
pure harmonic entropy). \\
\flg{--fscale <F>} & Vibrational frequency scale factor (default 1.0; must be
$>0$). \\
\flg{--makereal} & Treat imaginary frequencies as their absolute values
inside the derivation. \\
\flg{--ignoreimag} & Drop imaginary frequencies from the vibrational sums.
Conflicts with \flg{--makereal}. \\
\flg{--unit} & Print energies in kcal/mol instead of hartree. Takes no value. \\
\bottomrule
\end{longtable}

\subsection{Corrections and extrapolation}

\begin{longtable}{@{}L{4.3cm}L{9.2cm}@{}}
\toprule
\textbf{Flag} & \textbf{Meaning} \\
\midrule
\endhead
\flg{--add-correction <PATH>} & Attach single-point corrections from a
database, matching by exact case-sensitive \fld{name}. One path; not
repeatable. \\
\flg{--extrapolate <MODE VALUES>...} & Attach an extrapolated correction.
Modes: \texttt{2p LOW HIGH N ALPHA BETA}, \texttt{2s LOW HIGH F}, or
\texttt{2p LOWPAIR HIGHPAIR N ALPHA BETA F}. Place it after the positional
inputs. \\
\bottomrule
\end{longtable}

\subsection{Formation}

\begin{longtable}{@{}L{4.3cm}L{9.2cm}@{}}
\toprule
\textbf{Flag} & \textbf{Meaning} \\
\midrule
\endhead
\flg{--formation} & Print the raw formation-energy section. \\
\flg{--formation-corr} & Print the corrected formation-energy section. \\
\flg{--conc <CODE> <VALUE>} & Evaluate formation free energies at a vapor
concentration. Repeatable; one setting per code. \\
\flg{--cnt} & Use the self-consistent $(N-1)/N$ concentration scaling. \\
\flg{--bavg} & Collapse each printed formation section into one
Boltzmann-averaged row per composition. \\
\bottomrule
\end{longtable}

\subsection{Solvation and hydration}

\begin{longtable}{@{}L{4.3cm}L{9.2cm}@{}}
\toprule
\textbf{Flag} & \textbf{Meaning} \\
\midrule
\endhead
\flg{--hydration} & Print the raw water distribution section (solvent
\texttt{w}). \\
\flg{--hydration-corr} & Print the corrected water distribution section. \\
\flg{--solvation <CODE>} & Print the raw distribution section for solvent
\texttt{CODE}. \\
\flg{--solvation-corr <CODE>} & Print the corrected distribution section for
solvent \texttt{CODE}. \\
\flg{--rh <VALUE>} & Relative humidity as a fraction (0.5 is 50 \%). Water
only. \\
\flg{--psolvent <VALUE>} & Solvent partial pressure in the \flg{--conc}
VALUE grammar. Overrides \flg{--rh}; required for non-water solvents. \\
\bottomrule
\end{longtable}

\subsection{Help and version}

\begin{longtable}{@{}L{3.5cm}L{10.0cm}@{}}
\toprule
\textbf{Flag} & \textbf{Meaning} \\
\midrule
\endhead
\flg{-h}, \flg{--help} & Print help and exit. \\
\flg{-V}, \flg{--version} & Print the version and exit. \\
\bottomrule
\end{longtable}

\section{Database record reference}\label{app:record-reference}

A record is the information for one calculation, written as text on one line.
This appendix is the quick reference; Chapter~\ref{sec:jsonl} explains how to
read a database without programming.

\subsection{Entry order and presence}

\CJKQC{} writes the entries in this fixed order, leaving out information that
is not available:

\begin{verbatim}
status, error, name, source_path, program, program_version, command,
electronic_energy, scf_energy, charge, multiplicity, geometry,
vibrational_frequencies, gradient, dipole_moment, mulliken_charges,
rotational_symmetry_number, components, correction
\end{verbatim}

The entries may appear in any order, and unknown entries are ignored. Missing
information is left out completely; you will not see \texttt{null}. Entry
names use underscores and do not change between versions.

\subsection{Entries}

\begin{longtable}{@{}L{6.0cm}L{2.7cm}L{5.3cm}@{}}
\toprule
\textbf{Entry} & \textbf{Kind of value} & \textbf{Meaning} \\
\midrule
\endhead
\fld{status} & text & \fld{ok}, \fld{failed}, or \fld{truncated}. A missing
entry means \fld{ok}. \\
\fld{error} & text & Reason for a failure; never empty when
\fld{status} is \fld{failed}. \\
\fld{name} & text & File stem, exactly as written. The name used to match
corrections and monomer references. \\
\fld{source_path} & text & Full path of the original output file. \\
\fld{program} & text & Program name: \fld{gaussian}, \fld{orca}, \fld{xtb},
\fld{mrcc}, \fld{xyz}. \\
\fld{program_version} & text & Version as printed by the program. \\
\fld{command} & text & Command line as echoed by the program. \\
\fld{electronic_energy} & number & Final electronic energy, hartree. \\
\fld{scf_energy} & number & Converged SCF energy, hartree. Used only by
extrapolation. \\
\fld{charge} & whole number & Net charge in elementary charges. \\
\fld{multiplicity} & whole number & Spin multiplicity, at least 1. \\
\fld{geometry} & two lists & \fld{symbols} and \fld{coordinates}; coordinates
in \AA{}, listed atom by atom. \\
\fld{vibrational_frequencies} & list & Genuine modes in cm$^{-1}$; a negative
value means imaginary. \\
\fld{gradient} & list & Cartesian gradient, hartree/\AA{}, three values per
atom. \\
\fld{dipole_moment} & list & $[x,y,z]$ in atomic units ($e\,a_0$). \\
\fld{mulliken_charges} & list & Per-atom charges in elementary charges. \\
\fld{rotational_symmetry_number} & whole number & $\sigma$, at least 1. \\
\fld{components} & code and count & Each component code with its count, in
alphabetical order. \\
\fld{correction} & group of values & Attached correction; see below. \\
\bottomrule
\end{longtable}

\subsection{Geometry entries}

\begin{verbatim}
"geometry": {
  "symbols": ["N", "H"],
  "coordinates": [0.0, 0.0, 0.0, 0.0, 0.0, 1.01]
}
\end{verbatim}

\fld{symbols} holds the element names, with the first letter capital. The
\fld{coordinates} list has exactly three numbers per symbol, in the same atom
order: $x,y,z$ for the first atom, then $x,y,z$ for the second, and so on. The
number of coordinate values must be exactly three times the number of symbols.

\subsection{Components entry}

\begin{verbatim}
"components": {"am": 2, "sa": 3}
\end{verbatim}

Component codes are made of letters (A--Z and a--z), and upper and lower case
matter. Counts are at least 1. The codes are written in alphabetical order.

\subsection{Correction entries}

\begin{verbatim}
"correction": {
  "energy": -76.4109,
  "program": "orca",
  "program_version": "6.0.0",
  "commands": ["! wB97x-3c def2-TZVPP Sp"],
  "kind": "sp"
}
\end{verbatim}

\begin{longtable}{@{}L{4.2cm}L{2.3cm}L{7.0cm}@{}}
\toprule
\textbf{Entry} & \textbf{Kind of value} & \textbf{Meaning} \\
\midrule
\endhead
\fld{energy} & number & Correction energy in hartree. \\
\fld{program} & text & Program name of the correction calculation. \\
\fld{program_version} & text & Version of the correction calculation. \\
\fld{commands} & list & Command lines that produced the correction, in order. \\
\fld{kind} & text & How the energy was obtained: \texttt{sp}, \texttt{cbs},
\texttt{cps}, or \texttt{cbs-cps}. \\
\bottomrule
\end{longtable}

\subsection{When a record is rejected}

A record is rejected (and, during collection, replaced by a failed record) if
any of the following holds:

\begin{itemize}
  \item a number is not a finite value (it is infinitely large, or not a
        number at all);
  \item the geometry has no symbols, or its coordinate count is not exactly
        three times the symbol count;
  \item a per-atom list (\fld{gradient}, \fld{mulliken_charges}) disagrees
        with the geometry's atom count when a geometry is present;
  \item a \fld{failed} record has no error message;
  \item \fld{multiplicity} or \fld{rotational_symmetry_number} is below 1;
  \item \fld{components} is empty, a code is empty or contains non-letters, or
        a count is below 1;
  \item a correction has a non-finite energy or an empty \fld{kind}.
\end{itemize}

\subsection{Example records}

A minimal record (the actual file stores this on one line):

\begin{verbatim}
{"status":"ok","electronic_energy":-1.5}
\end{verbatim}

A record with a geometry and an attached correction, shown spread over several
lines for readability:

\begin{verbatim}
{
  "status": "ok",
  "name": "1dma",
  "program": "xtb",
  "electronic_energy": -38.51,
  "geometry": {
    "symbols": ["N", "H"],
    "coordinates": [0.0, 0.0, 0.0, 0.0, 0.0, 1.01]
  },
  "components": {"dma": 1},
  "correction": {
    "energy": -76.4109,
    "program": "orca",
    "program_version": "6.0.0",
    "commands": ["! wB97x-3c def2-TZVPP Sp"],
    "kind": "sp"
  }
}
\end{verbatim}

\section{Constants, exit codes, and messages}\label{app:constants}

\subsection{Physical constants}

All conversions use the CODATA 2018 values below. Energies are held in
hartree internally and converted only at the print boundary.

\begin{center}
\begin{tabular}{@{}ll@{}}
\toprule
\textbf{Constant} & \textbf{Value} \\
\midrule
Planck constant $h$ & $6.62607015\times10^{-34}$ J s \\
Speed of light $c$ & $2.99792458\times10^{10}$ cm s$^{-1}$ \\
Boltzmann constant $k_B$ & $1.380649\times10^{-23}$ J K$^{-1}$ \\
Avogadro constant $N_A$ & $6.02214076\times10^{23}$ mol$^{-1}$ \\
Gas constant $R = N_A k_B$ & $8.31446261815324$ J mol$^{-1}$ K$^{-1}$ \\
Atomic mass constant $u$ & $1.66053906660\times10^{-27}$ kg \\
Standard pressure $p^\circ$ (1 atm) & 101325 Pa \\
Hartree energy $E_h$ & $4.3597447222071\times10^{-18}$ J \\
1 kcal/mol & 4184 J/mol (exact) \\
Hartree to kcal/mol & 627.5094740631 \\
Bohr to \AA{} & 0.529177210903 \\
Debye per $e\,a_0$ & 2.541746473 \\
\bottomrule
\end{tabular}
\end{center}

\subsection{Exit codes}

\begin{center}
\begin{tabular}{@{}cl@{}}
\toprule
\textbf{Code} & \textbf{Meaning} \\
\midrule
0 & The run completed and skipped nothing. \\
1 & A fatal error occurred and nothing was written, or a merge skipped
corrupt lines. \\
2 & Usage error: bad flags, missing \flg{--out}, mixed input kinds, or an
invalid value. \\
\bottomrule
\end{tabular}
\end{center}

Failed and truncated records from a recognized format do not change the exit
code. Warnings about unmatched correction candidates also do not change it.

\subsection{Normal output messages}

\begin{longtable}{@{}L{6.2cm}L{7.2cm}@{}}
\toprule
\textbf{Message} & \textbf{When} \\
\midrule
\endhead
\texttt{collected <n> records into <out> (<f> failed, <t> truncated)} &
A build run completed. \\
\texttt{collected <n> records into <out> (<f> failed, <t> truncated, <k>
corrected)} & A build run with \flg{--add-correction} completed. \\
\texttt{merged <n> records into <out>} & A merge run completed. \\
\texttt{merged <n> records into <out> (<k> corrected)} & A merge run with
corrections completed. \\
\texttt{filtered <n> records into <out>} & A manipulation run wrote records. \\
\texttt{filtered <n> records into <out> (<k> corrected)} & A manipulation run
with corrections wrote records. \\
\texttt{wrote <n> frames into movie.xyz} & \flg{--movie} wrote a file. \\
\bottomrule
\end{longtable}

A table, formation section, corrected formation section, and
hydration/solvation sections print before the completion line.

\subsection{Warning and error messages}

\begin{longtable}{@{}L{7.0cm}L{6.4cm}@{}}
\toprule
\textbf{Message} & \textbf{Meaning} \\
\midrule
\endhead
\texttt{excluded <k> records with status failed or truncated} & The filtering
steps dropped non-\fld{ok} records before any other step. \\
\texttt{warning: <path> line <n>: <reason>} & A corrupt database line was
skipped. \\
\texttt{warning: ... correction ...} & A correction candidate was skipped or
unmatched. Warnings do not fail the run. \\
\texttt{<path> already exists; overwrite? [y/N]} & Interactive overwrite
confirmation. \\
\texttt{error: <path> exists; not overwriting} & The overwrite answer was not
yes. \\
\texttt{error: <path> exists; not interactive; use -{}-fout to force} & The
the run is not interactive (input is not a terminal). \\
\texttt{excluded <n> records from movie.xyz (no geometry or not ok)} & Some
records could not become frames. \\
\bottomrule
\end{longtable}

\subsection{Common fatal error texts}

The exact wording can evolve, but these are the recurring causes:

\begin{itemize}
  \item an input could not be opened (missing file, directory, permission
        denied);
  \item an input matched no supported format, or was empty;
  \item \flg{--out} could not be written, or its parent directory does not
        exist;
  \item a correction or extrapolation source was missing, unreadable, or
        ambiguous;
  \item a correction database or source loaded no record while reporting
        load problems;
  \item two extrapolation candidates shared a name, mixed programs or
        versions, or repeated an electronic energy exactly.
\end{itemize}

\section{Troubleshooting}\label{app:troubleshooting}

\subsection{``Mixed inputs'' or exit code 2}

An invocation is either all program outputs (build) or all databases (merge).
Mixing the two kinds is a usage error. Run the two steps separately:

\begin{verbatim}
CJKQC *.out -o part.jsonl --fout
CJKQC part.jsonl other.jsonl -o all.jsonl --fout
\end{verbatim}

Classification is by content, not extension, so a file whose first non-blank
line begins with \texttt{\{} counts as a database.

\subsection{An input is not recognized}

A build input that matches no supported format, or is empty, aborts the run
and writes nothing. Check:

\begin{itemize}
  \item that the file is really a program output and not, for example, a
        shell script or a truncated header;
  \item that the file is not empty;
  \item that a wrapper has not hidden the program banner beyond the head
        window; try increasing \flg{--head-bytes}.
\end{itemize}

A failed or truncated \emph{job} is different: if \CJKQC{} recognizes the
format, the record is stored and reported in the summary line.

\subsection{``cut by the window seam''}

A geometry or frequency table was split between the head and tail windows.
Widen the windows and retry:

\begin{verbatim}
CJKQC huge.out --head-bytes 1048576 --tail-bytes 8388608 -o db.jsonl --fout
\end{verbatim}

The defaults cover the shipped fixture corpus, but unusually large or
unusually laid-out outputs may need more.

\subsection{Thermochemistry prints dashes}

An undefined derivation is not an error. Check for:

\begin{itemize}
  \item a missing \fld{electronic_energy};
  \item a missing geometry, or a degenerate geometry;
  \item an unknown element symbol in the geometry;
  \item a negative frequency without \flg{--makereal} or
        \flg{--ignoreimag};
  \item a non-monatomic record with no participating vibrational mode, such
        as a single-point output or a plain xyz file.
\end{itemize}

Try \flg{--makereal} to treat an imaginary mode as real, or
\flg{--ignoreimag} to drop it. Only one may be used at a time.

\subsection{``\texttt{-{}-pop} db.jsonl fails''}

\flg{--pop} takes an optional level, and a space-separated value is consumed
even if it is your database path. Put \flg{--pop} after the positional inputs,
or use the equals form:

\begin{verbatim}
CJKQC db.jsonl --pop
CJKQC --pop=g db.jsonl
\end{verbatim}

The same value-eating behavior applies to \flg{--extrapolate}; place it after
the positional inputs.

\subsection{A correction does not attach}

Matching is exact and case-sensitive on \fld{name}. Common causes:

\begin{itemize}
  \item the correction record's name differs in case or spelling;
  \item the run record has no \fld{name};
  \item the correction record is not \fld{ok}, or lacks
        \fld{electronic_energy}, so it is not a candidate;
  \item the correction record's name is absent.
\end{itemize}

Read the standard-error warnings: skipped candidates and unmatched candidates
are reported individually or as an aggregate count.

\subsection{Extrapolation refuses to run}

Extrapolation rejects ambiguity rather than producing a silent mix:

\begin{itemize}
  \item two candidates with the same name in one source database;
  \item contributors of one name with different programs or versions;
  \item contributors of one name with exactly equal electronic energies,
        which usually means the same file entered two sources or the same
        path was passed twice;
  \item a run record whose name is missing from at least one source. A
        partially covered name is warned and not corrected.
\end{itemize}

Every source path must be a database. Use \flg{--add-correction} for a single
ordinary correction; a run cannot use both.

\subsection{``exists; not interactive; use \texttt{-{}-fout} to force''}

An output file already exists and standard input is not a terminal, so
\CJKQC{} will not ask. Either run interactively and answer \texttt{y}, or add
\flg{--fout}. The same rule applies to \texttt{movie.xyz}. A refusal at
\flg{--out} ends the run before \texttt{movie.xyz} is considered.

\subsection{A formation block or solvation block is missing}

Formation and solvation need monomers in the run's final record set:

\begin{itemize}
  \item a monomer is an \fld{ok} record with exactly one component at count 1;
  \item if a filter removed the monomer, it cannot serve as a reference;
  \item for solvation, the solvent code must appear in at least one record of
        the section's record set;
  \item a state whose solvent monomer is absent or has no level value prints a
        dash.
\end{itemize}

For non-water solvents, supply \flg{--psolvent}; the water saturation
correlation is not applied to other solvents.

\subsection{Sorting by name does not work}

Use \flg{--sort filename}. The accepted value is \fld{filename}, which sorts
by the record's \fld{name}. The help text's phrase ``name'' refers to that
name, not to a separate spelling.

\subsection{Populations or formation values changed after filtering}

Populations, formation references, and solvation distributions are computed
over exactly the records the run prints. Adding a filter changes the set, so
the normalization and the selected reference monomers can change too. This is
intended: these are ensemble properties of the printed set, not stored
per-record values.


\end{document}
